% Programs as Histories: A Case Study in Ontology-Driven Language Design
% Flyxion, June 2026
% Compiled with LuaLaTeX

\documentclass[12pt,oneside,openany]{book}

%% ── Fonts ────────────────────────────────────────────────────────────────────
\usepackage{fontspec}
\setmainfont{TeX Gyre Pagella}
\usepackage{booktabs}
\usepackage{unicode-math}
\setmathfont{Latin Modern Math}

%% ── Layout ───────────────────────────────────────────────────────────────────
\usepackage[
  paperwidth=6.5in, paperheight=9.5in,
  top=1.1in, bottom=1.1in,
  inner=1in, outer=0.85in,
  marginparwidth=0pt
]{geometry}
\usepackage{microtype}
\frenchspacing
\setlength{\parskip}{0pt}
\setlength{\parindent}{1.4em}

%% ── Mathematics ──────────────────────────────────────────────────────────────
\usepackage{amsmath}
\usepackage{amsthm}
\usepackage{mathtools}

%% ── Theorem environments ─────────────────────────────────────────────────────
\newtheoremstyle{spherepop}
  {8pt}{6pt}{\itshape}{}{\bfseries}{.}{0.5em}{}
\theoremstyle{spherepop}
\newtheorem{definition}{Definition}[chapter]
\newtheorem{theorem}[definition]{Theorem}
\newtheorem{proposition}[definition]{Proposition}
\newtheorem{lemma}[definition]{Lemma}
\newtheorem{corollary}[definition]{Corollary}

\newtheoremstyle{remarkstyle}
  {6pt}{4pt}{\upshape}{}{\bfseries}{.}{0.5em}{}
\theoremstyle{remarkstyle}
\newtheorem{remark}[definition]{Remark}
\newtheorem{example}[definition]{Example}
\newtheorem{notation}[definition]{Notation}

%% ── Code listings ────────────────────────────────────────────────────────────
\usepackage{listings}
\usepackage{xcolor}

\definecolor{codebg}{RGB}{248,248,246}
\definecolor{codekw}{RGB}{140,40,90}
\definecolor{codecomment}{RGB}{100,110,100}
\definecolor{codestring}{RGB}{60,100,160}
\definecolor{codenumber}{RGB}{90,90,90}

\lstdefinelanguage{Rust}{
  keywords={pub,fn,let,mut,match,enum,struct,impl,use,mod,
            where,type,trait,for,in,if,else,return,Ok,Err,
            Box,Vec,Option,Some,None,Result,Self,self,super,crate},
  keywordstyle=\color{codekw}\bfseries,
  comment=[l]{//},
  commentstyle=\color{codecomment}\itshape,
  stringstyle=\color{codestring},
  morestring=[b]",
  sensitive=true,
}

\lstset{
  language=Rust,
  basicstyle=\small\ttfamily,
  backgroundcolor=\color{codebg},
  frame=none,
  xleftmargin=1.5em,
  xrightmargin=0.5em,
  aboveskip=8pt,
  belowskip=6pt,
  breaklines=true,
  showstringspaces=false,
  numbers=none,
  tabsize=4,
}

%% ── Section formatting ───────────────────────────────────────────────────────
\usepackage{titlesec}
\titleformat{\chapter}[display]
  {\large\bfseries\raggedright}
  {\normalsize\mdseries\scshape\chaptertitlename\ \thechapter}
  {6pt}
  {\LARGE\bfseries}
  [\vspace{4pt}\titlerule]
\titlespacing*{\chapter}{0pt}{24pt}{20pt}

\titleformat{\section}{\large\bfseries}{\thesection}{0.8em}{}
\titleformat{\subsection}{\normalsize\bfseries\itshape}{\thesubsection}{0.7em}{}

%% ── Headings, TOC ────────────────────────────────────────────────────────────
\usepackage{fancyhdr}
\pagestyle{fancy}
\fancyhf{}
\fancyhead[LE]{\small\itshape\leftmark}
\fancyhead[RO]{\small\itshape\rightmark}
\fancyfoot[C]{\small\thepage}
\renewcommand{\headrulewidth}{0.3pt}

\usepackage{tocloft}
\setlength{\cftbeforechapskip}{4pt}

%% ── Epigraph ─────────────────────────────────────────────────────────────────
\usepackage{epigraph}
\setlength{\epigraphwidth}{0.7\textwidth}
\renewcommand{\epigraphflush}{center}
\renewcommand{\epigraphrule}{0pt}

%% ── Hyperref (last) ──────────────────────────────────────────────────────────
\usepackage[
  colorlinks=true,
  linkcolor=black,
  citecolor=black,
  urlcolor=black,
  pdftitle={Programs as Histories},
  pdfauthor={Flyxion},
]{hyperref}

%% ── Custom operators ─────────────────────────────────────────────────────────
\newcommand{\Hist}{\mathbf{Hist}}
\newcommand{\Obs}{\mathbf{Obs}}
\newcommand{\Pop}{\mathrm{Pop}}
\newcommand{\Rfuse}{\mathrm{Refuse}}
\newcommand{\Bind}{\mathrm{Bind}}
\newcommand{\Col}{\mathrm{Collapse}}
\newcommand{\hconcat}{\mathbin{+\mkern-5mu+}}  % history concat
\newcommand{\Adm}[1]{\mathrm{Admissible}(#1)}
\newcommand{\Refused}[2]{\mathrm{Refused}(#1,\,#2)}
\newcommand{\Collapsed}[2]{\mathrm{Collapsed}(#1,\,#2)}
\newcommand{\vdashA}{\vdash_{\!\mathcal{A}}}
\newcommand{\yields}{\leadsto}
\newcommand{\calA}{\mathcal{A}}
\newcommand{\calH}{\mathcal{H}}
\newcommand{\calI}{\mathcal{I}}

%% ── Tcolorbox for implementation boxes ───────────────────────────────────────
\usepackage{tcolorbox}
\tcbuselibrary{skins,breakable}
\newtcolorbox{implbox}[1]{
  colback=codebg,
  colframe=black!20,
  fonttitle=\small\bfseries,
  title={#1},
  arc=2pt,
  left=6pt, right=6pt, top=4pt, bottom=4pt,
  breakable,
}

% ─────────────────────────────────────────────────────────────────────────────
\begin{document}

\frontmatter

%% ── Title ────────────────────────────────────────────────────────────────────
\begin{titlepage}
\vspace*{2cm}
\begin{center}
{\fontsize{26}{32}\selectfont\bfseries Programs as Histories}\\[6pt]
{\fontsize{18}{24}\selectfont\bfseries\itshape Projection, Admissibility, and the Geometry of Computation}\\[12pt]
{\large A Case Study in Ontology-Driven Language Design}\\[48pt]
{\large Flyxion}\\[8pt]
{\normalsize Independent Research}\\[48pt]
{\normalsize June 2026}
\end{center}
\vfill
\begin{center}
\rule{0.5\textwidth}{0.4pt}\\[6pt]
{\small\itshape
The implementation did not merely realize the theory;\\
it exposed which distinctions the theory was still hiding.}
\end{center}
\end{titlepage}

%% ── Abstract ─────────────────────────────────────────────────────────────────
\chapter*{Abstract}
\addcontentsline{toc}{chapter}{Abstract}

Most programming languages begin with values, variables, and state transitions.
Spherepop begins elsewhere.
It begins with the claim that histories are ontologically prior to states,
and that computation should be understood as the progressive restriction of
possibility rather than the manipulation of values.

This monograph presents the design and formal theory of the Spherepop
programming language and interpreter.
We proceed from ontological commitment through mathematical formalization
to implementation, treating each software component as the computational
expression of a prior philosophical necessity.

We establish a formal operational semantics, a typed admissibility judgement
system, four principal theorems (History Monotonicity, Option Space
Monotonicity, Collapse Soundness, Replay Equivalence), and a
category-theoretic account of histories as morphisms and collapse as
a functor between history and observable-state categories.

Four implementation discoveries retroactively sharpened the theory:
refusal must carry reasons; collapse must specify quotient rules;
type checking depends on the assumed boundary of the world;
and compiler correctness requires history equivalence, not mere output equivalence.

The thesis is outrageous only as a slogan.
As a design principle it is precise:
\emph{computation is the progressive restriction of possibility,
not the manipulation of values.}

\tableofcontents

\mainmatter

% ═══════════════════════════════════════════════════════════════════════════════
\part{Why State Is the Wrong Primitive}
% ═══════════════════════════════════════════════════════════════════════════════

\chapter{The State-Centric Assumption}

\epigraph{\itshape
A program is a sequence of state transformations.
So runs the unexamined premise of nearly every language ever designed.}{---}

\section{Motivation}

Open any introductory programming text and you encounter a common first lesson:
a variable is a named location that holds a value.
A program begins in some initial state,
instructions modify that state,
and the final state is the answer.

\begin{lstlisting}
x = 5
y = x + 1
\end{lstlisting}

This appears natural.
It mirrors how we naively think about physical processes:
an object has properties; those properties change;
the current properties constitute the object's state.
But the appearance of naturalness is produced by a prior choice ---
the choice to treat present configuration as the fundamental unit of description.

\section{The State Illusion}

Consider what this choice erases.
In version control, the repository is not merely its current file tree;
it is the entire sequence of commits that produced it.
In event sourcing, database records are append-only event logs;
the current state is a fold over that log.
In scientific experiment, a measurement is the terminus of a sequence
of decisions, calibrations, and interpretive choices.
In evolutionary biology, a genome is a compressed record of selective pressures.

In each case, the historical structure is epistemically and causally prior
to the present configuration.
State is the summary.
History is the substance.

\begin{definition}[State Degeneracy]
\label{def:degeneracy}
Let $E$ be an event alphabet, $\calH = \bigcup_{n=0}^{\infty} E^n$
the set of all finite histories, and $S$ an observable state space.
A \emph{state projection} is a function $\sigma : \calH \to S$.
The \emph{degeneracy} of a state $s \in S$ is
\[
D(s) \;=\; \bigl|\sigma^{-1}(s)\bigr|.
\]
\end{definition}

\begin{proposition}[The State Illusion]
\label{prop:state-illusion}
For any nontrivial finite state space $S$ and any non-injective projection
$\sigma : \calH \to S$, there exist distinct histories $H_1 \neq H_2$ such that
$\sigma(H_1) = \sigma(H_2)$.
Consequently,
\[
\sigma(H_1) = \sigma(H_2) \;\not\Rightarrow\; H_1 = H_2.
\]
\end{proposition}

\begin{proof}
Since $|\calH|$ is countably infinite and $|S|$ is finite, $\sigma$ cannot be injective.
Therefore $D(s) > 1$ for at least one $s \in S$,
which yields the required $H_1 \neq H_2$ with $\sigma(H_1) = \sigma(H_2)$.
\end{proof}

\begin{remark}
Proposition~\ref{prop:state-illusion} is the formal statement of what we call
the \emph{state illusion}: any program that exposes only $\sigma(H)$ to its
user has silently discarded distinguishing information about the history $H$.
The conventional debugger is an instrument for reconstructing $H$ from $\sigma(H)$
after information has already been lost.
\end{remark}

\section{Consequences for Language Design}

Programming languages that treat state as primary inherit the degeneracy
of Proposition~\ref{prop:state-illusion} as a structural feature.
Race conditions arise because $\sigma$ is non-injective on interleaved histories.
Use-after-free vulnerabilities arise because $\sigma$ conflates histories
that differ only in deallocation events.
The auditing and provenance tools that practitioners add to production
systems are, in each case, partial reconstructions of $H$ from $\sigma(H)$.

The Spherepop design decision follows directly:
if history is what is lost, history must be what is preserved.

% ─────────────────────────────────────────────────────────────────────────────
\chapter{The Spherepop Inversion}

\epigraph{\itshape
State is a collapsed history.
The collapse is the problem.}{---}

\section{Motivation}

Spherepop proposes the opposite ontological priority from conventional languages.
Not $\text{State} \to \text{History}$ (history as debugging afterthought)
but $\text{History} \to \text{State}$ (state as derived projection).

\section{The World Structure}

\begin{definition}[Event Alphabet]
Let $E$ be a set of \emph{events}, partitioned into primitive generators:
\[
E = \{\Pop\} \cup \{\Rfuse\} \cup \{\Bind\} \cup \{\Col\} \cup E_\lambda
\]
where $E_\lambda$ contains the lambda-calculus events
$\{\mathrm{LamIntro}, \mathrm{Apply}\}$ and scope markers.
\end{definition}


\begin{notation}
Throughout, $\Omega_0$ denotes the \emph{initial finite option space} --- the
full set of possibilities available at the start of a computation.
All option spaces $\Omega$ appearing in subsequent definitions satisfy
$\Omega \subseteq \Omega_0$ and $|\Omega_0| < \infty$.
\end{notation}

\begin{definition}[Spherepop World]
\label{def:world}
A \emph{Spherepop world} is a pair
\[
W \;=\; (H,\,\Omega)
\]
where $H \in \calH$ is a history (finite event sequence) and
$\Omega \subseteq \Omega_0$ is the current \emph{option space}
(a finite set of available symbols), derived from some initial option space $\Omega_0$.
\end{definition}

\begin{definition}[Commitment Operator]
\label{def:pop-op}
The \emph{commitment operator} for symbol $x \in \Omega$ is
\[
P_x : (H,\Omega) \;\mapsto\; \bigl(H \hconcat [\Pop(x)],\; \Omega \setminus \{x\}\bigr)
\]
where $\oplus$ denotes history concatenation.
$P_x$ is undefined when $x \notin \Omega$.
\end{definition}

\begin{definition}[Refusal Operator]
\label{def:refuse-op}
The \emph{refusal operator} for symbol $x$ with reason $r$ is
\[
R_{x,r} : (H,\Omega) \;\mapsto\; \bigl(H \hconcat [\Rfuse(x,r)],\; \Omega\bigr).
\]
Note that $\Omega$ is \emph{unchanged}: refusal records inadmissibility
without foreclosing structural availability.
\end{definition}

\begin{definition}[Bind and Collapse Operators]
\[
B_{a,b} : (H,\Omega) \;\mapsto\; (H \hconcat [\Bind(a,b)],\; \Omega)
\]
\[
C_{x,c} : (H,\Omega) \;\mapsto\; (H \hconcat [\Col(x,c)],\; \Omega)
\]
where $c$ is a collapse rule (see Definition~\ref{def:collapse-rule}).
\end{definition}

\section{The Conservation Law}

\begin{theorem}[Conservation of Possibility]
\label{thm:conservation}
For any legal execution sequence
$W_0 \xrightarrow{e_1} W_1 \xrightarrow{e_2} \cdots \xrightarrow{e_n} W_n$
where each step is one of $P_x$, $R_{x,r}$, $B_{a,b}$, or $C_{x,c}$:
\begin{enumerate}
\item[(i)] $|H_t|$ is strictly monotonically increasing: $|H_{t+1}| = |H_t| + 1$.
\item[(ii)] $|\Omega_t|$ is monotonically non-increasing.
\item[(iii)] Under pure Pop dynamics, $|H_t| + |\Omega_t| = |\Omega_0|$ for all $t$.
\end{enumerate}
\end{theorem}

\begin{proof}
(i) Every operator appends exactly one event to $H$, so $|H_{t+1}| = |H_t| + 1$.

(ii) $P_x$ removes $x$ from $\Omega$; the remaining operators do not modify $\Omega$.
Hence $|\Omega_{t+1}| \leq |\Omega_t|$.

(iii) Under pure Pop dynamics, every step removes exactly one element from $\Omega$
and appends one event to $H$.
Starting from $|H_0| = 0$ and $|\Omega_0|$,
after $t$ steps $|H_t| = t$ and $|\Omega_t| = |\Omega_0| - t$,
giving $|H_t| + |\Omega_t| = |\Omega_0|$.
\end{proof}

\begin{remark}
Theorem~\ref{thm:conservation} is the Spherepop analogue of a conservation law
in physics. It says that possibility is not created or destroyed ---
it is either consumed (by Pop) or documented as inadmissible (by Refuse).
The total ``computational material'' $|H_t| + |\Omega_t|$ is invariant
under pure commitment dynamics.
\end{remark}

\begin{implbox}{Implementation consequence: \texttt{core/world.rs}}
The conservation law forces the World to own both \texttt{history} and
\texttt{options} as co-equal fields. Observable state is absent from the
struct --- it is always derived on demand.
\begin{lstlisting}
pub struct World {
    pub options: OptionSpace,
    pub history: History,
}
\end{lstlisting}
The absence of a \texttt{state} field is the mechanization of
Theorem~\ref{thm:conservation}: observable state is never fundamental.
\end{implbox}

% ═══════════════════════════════════════════════════════════════════════════════
\part{Histories Before Computation}
% ═══════════════════════════════════════════════════════════════════════════════

\chapter{The Algebra of Histories}

\section{Motivation}

Before introducing types, evaluators, or compilers, we must characterize
the mathematical structure of the objects that are primary in our ontology.
Histories are not merely lists of events.
They form a category, and the category structure constrains what operations
on histories are meaningful.

\section{Definitions}

\begin{definition}[History]
A \emph{history} over event alphabet $E$ is a finite sequence
$H = [e_1, e_2, \ldots, e_n]$ with each $e_i \in E$.
The empty history is denoted $\varepsilon = []$.
\end{definition}

\begin{definition}[History Concatenation]
For histories $H_1 = [e_1, \ldots, e_m]$ and $H_2 = [f_1, \ldots, f_n]$,
\[
H_1 \hconcat H_2 = [e_1, \ldots, e_m, f_1, \ldots, f_n].
\]
\end{definition}

\begin{definition}[History Monoid]
The triple $(\calH, \hconcat, \varepsilon)$ where $\calH = \bigcup_{n \geq 0} E^n$
is the \emph{history monoid}.
\end{definition}

\begin{proposition}
$(\calH, \hconcat, \varepsilon)$ is a free monoid over $E$.
\end{proposition}

\begin{proof}
Associativity of $\oplus$ follows from list concatenation associativity.
$\varepsilon$ is the left and right identity.
Freeness: every history $H$ is uniquely determined by its event sequence,
and no non-trivial relation among generators holds.
\end{proof}

\section{The History Category}

\begin{definition}[History Category $\Hist$]
\label{def:hist-category}
Define the category $\Hist$ as follows:
\begin{itemize}
\item \emph{Objects:} option spaces $\Omega \subseteq \Omega_0$.
\item \emph{Morphisms:} $\mathrm{Hom}(\Omega, \Omega')$ is the set of histories
  $H$ such that starting from option space $\Omega$, executing $H$ yields $\Omega'$.
\item \emph{Composition:} $H_2 \circ H_1 = H_1 \hconcat H_2$
  (execute $H_1$ first, then $H_2$).
\item \emph{Identity:} $\mathrm{id}_\Omega = \varepsilon$.
\end{itemize}
\end{definition}

\begin{proposition}
$\Hist$ is a well-defined category.
\end{proposition}

\begin{proof}
Associativity of composition follows from associativity of $\oplus$.
$\varepsilon$ is the identity morphism since $H \hconcat \varepsilon = H = \varepsilon \hconcat H$.
\end{proof}

\begin{remark}
$\Hist$ is in fact a \emph{free strict monoidal category}:
tensor product of morphisms is history concatenation (Meld),
and the monoidal unit is the empty option space with empty history.
This is much stronger than the simple monoid observation:
it equips histories with a compositional structure that supports
parallel as well as sequential combination.
\end{remark}

\begin{implbox}{Implementation consequence: \texttt{core/history.rs}}
The free monoid structure forces \texttt{append} (extend by one generator)
and \texttt{meld} (monoidal composition of two histories) as the only
legitimate history operations. The absence of \texttt{remove} and \texttt{undo}
is the mechanization of freeness: there are no relations among generators.
\begin{lstlisting}
pub fn append(&mut self, event: Event) {
    self.events.push(event);    // extend by one generator
}
pub fn meld(&mut self, other: &History) {
    self.events.extend(other.events.iter().cloned()); // monoidal +
}
\end{lstlisting}
\emph{History equality} $H_1 = H_2$ is the primary criterion of
program identity. This is implemented as structural list equality.
\end{implbox}

% ─────────────────────────────────────────────────────────────────────────────
\chapter{Operational Semantics}

\section{Motivation}

A formal operational semantics transforms the philosophical commitments of
Chapter~2 into mathematically precise transition rules.
Every execution step is a world-state transition;
the full execution is a path in the transition system.

\section{Small-Step Transition Rules}

We write $(H, \Omega) \xrightarrow{\alpha} (H', \Omega')$ for a single step
labelled $\alpha$.

\medskip
\noindent\textbf{Pop} (commitment):
\[
\frac{x \in \Omega}
{(H,\,\Omega)
  \;\xrightarrow{\mathrm{pop}(x)}\;
  (H \hconcat [\Pop(x)],\; \Omega \setminus \{x\})}
\quad \textsc{[Pop]}
\]

\noindent\textbf{Refuse} (documented inadmissibility):
\[
\frac{}
{(H,\,\Omega)
  \;\xrightarrow{\mathrm{refuse}(x,r)}\;
  (H \hconcat [\Rfuse(x,r)],\; \Omega)}
\quad \textsc{[Refuse]}
\]
Note: \textsc{[Refuse]} has no premise --- any symbol may be refused.
The effect on the option space is nil; the effect on admissibility is non-nil
(see Definition~\ref{def:admissibility-contraction}).

\noindent\textbf{Bind} (dependency declaration):
\[
\frac{}
{(H,\,\Omega)
  \;\xrightarrow{\mathrm{bind}(a,b)}\;
  (H \hconcat [\Bind(a,b)],\; \Omega)}
\quad \textsc{[Bind]}
\]

\noindent\textbf{Collapse} (observation under rule $c$):
\[
\frac{\Gamma \vdash t : \Adm{T}}
{(H,\,\Omega)
  \;\xrightarrow{\mathrm{collapse}(x,c)}\;
  (H \hconcat [\Col(x,c)],\; \Omega)}
\quad \textsc{[Collapse]}
\]
The premise $\Gamma \vdash t : \Adm{T}$ is the type-theoretic admissibility
guard (Chapter~6). Collapse requires an admissibility certificate.

\section{Admissibility as a Set-Valued Function}

\begin{definition}[Admissible Futures]
\label{def:admissible-futures}
For a history $H$, the set of \emph{admissible future symbols} is
\[
\calA(H) = \Omega_0 \setminus \bigl\{
  x \;\big|\; \exists r.\, \Rfuse(x,r) \in H
\bigr\}.
\]
\end{definition}

\begin{definition}[Inadmissibility Indicator]
\label{def:inadmissibility}
\[
\calI(H,x) = 1 - \chi_{\calA(H)}(x)
\]
where $\chi_A$ is the indicator function of set $A$.
\end{definition}

\begin{definition}[Admissibility Contraction]
\label{def:admissibility-contraction}
Let $H' = H \hconcat [\Rfuse(x,r)]$.
Then $\calA(H') \subseteq \calA(H)$,
and specifically $x \notin \calA(H')$ while $x$ may have been in $\calA(H)$.
\end{definition}

\begin{remark}
The asymmetry between \textsc{[Pop]} and \textsc{[Refuse]} is critical:
Pop contracts the \emph{structural} option space $\Omega$.
Refuse contracts the \emph{admissibility} set $\calA(H)$
without touching $\Omega$.
Two paths through possibility space may have the same $\Omega$ while
having different $\calA(H)$ --- they are structurally equivalent
but semantically distinguished.
\end{remark}

% ═══════════════════════════════════════════════════════════════════════════════
\part{Types as Admissibility Certificates}
% ═══════════════════════════════════════════════════════════════════════════════

\chapter{From Sets to Reachability Certificates}

\section{Motivation}

The standard interpretation of types --- a type is a set of admissible values ---
is adequate for state-centric languages.
It is insufficient for Spherepop.

Consider what a type must certify in a history-primary setting.
A value $v$ of type $T$ is not merely a member of a set.
It is the result of a path through history that reached $v$ without
violating any constraint represented by $T$.
The type is not a description of $v$ at a moment;
it is a certificate about the history that produced $v$.

\section{Definitions}

\begin{definition}[Spherepop Types]
\label{def:types}
The type language of Spherepop is defined inductively:
\begin{align*}
T ::= \;& \mathrm{Unit}
    \;\mid\; \mathrm{Never}
    \;\mid\; X \quad (X \text{ a type variable}) \\
    \;\mid\;& \Pi(x : T_1).T_2 \quad (\text{dependent function type}) \\
    \;\mid\;& \Adm{T} \quad (\text{admissibility certificate}) \\
    \;\mid\;& \Refused{T}{r} \quad (\text{documented inadmissibility}) \\
    \;\mid\;& \Collapsed{T}{c} \quad (\text{sealed observation under rule } c) \\
    \;\mid\;& \mathrm{Process}(T_1 \rightsquigarrow T_2)
\end{align*}
The simple function type $T_1 \to T_2$ is notation for $\Pi(\_ : T_1).T_2$.
\end{definition}

\begin{remark}
The distinctive types are $\Adm{T}$, $\Refused{T}{r}$, and $\Collapsed{T}{c}$.
They carry not just the inner type but a \emph{certificate}:
\begin{itemize}
\item $\Adm{T}$: the path that produced this term was admissible.
\item $\Refused{T}{r}$: the term was encountered but found inadmissible for reason $r$.
  The reason is part of the \emph{type} --- two terms refused for different reasons
  have genuinely different types.
\item $\Collapsed{T}{c}$: an admissible term was observed under collapse rule $c$
  and the observation is sealed.
  The rule $c$ is part of the type.
\end{itemize}
\end{remark}

% ─────────────────────────────────────────────────────────────────────────────
\chapter{The Admissibility Type System}

\section{Typing Judgements}

We write $\Gamma \vdash t : T$ for the standard typing judgement
(term $t$ has type $T$ in context $\Gamma$) and
$\Gamma \vdashA t : T$ for the \emph{admissibility judgement}
(the path to $t$ was admissible under $\Gamma$).

\section{Typing Rules}

\subsection{Pop: Issuing an Admissibility Certificate}

\[
\frac{\Gamma \vdash t : T}
{\Gamma \vdash \mathrm{pop}(t) : \Adm{T}}
\quad \textsc{[T-Pop]}
\]

Pop transforms any typed term into an admissible term.
It is the mechanism by which structural commitment produces a reachability certificate.

\subsection{Refuse: Recording Inadmissibility with Reason}

\[
\frac{\Gamma \vdash t : T \quad r \in \mathsf{RefusalReason}}
{\Gamma \vdash \mathrm{refuse}(t,r) : \Refused{T}{r}}
\quad \textsc{[T-Refuse]}
\]

Refusal is typed: the type carries the reason.
A term refused for \textsc{ConstraintViolation}$(c)$ has a type that is
\emph{formally distinct} from one refused for \textsc{NotAvailable}.

\subsection{Collapse: Sealing an Admissible Observation}

\[
\frac{\Gamma \vdash t : \Adm{T} \quad c \in \mathsf{CollapseRule}}
{\Gamma \vdash \mathrm{collapse}(t,c) : \Collapsed{T}{c}}
\quad \textsc{[T-Collapse]}
\]

Collapse requires an admissibility certificate.
A term whose type is not $\Adm{\cdot}$ cannot be collapsed --- this is a hard
type error, not a runtime exception.

\subsection{Bind: Recording a Dependency}

\[
\frac{\Gamma \vdash a : T_1 \quad \Gamma \vdash b : T_2}
{\Gamma \vdash \mathrm{bind}(a,b) : \mathrm{Process}(T_1 \rightsquigarrow T_2)}
\quad \textsc{[T-Bind]}
\]

\subsection{Lambda and Application}

\[
\frac{\Gamma, x : T_1 \vdash b : T_2}
{\Gamma \vdash \lambda x : T_1.b
  : \Pi(x : T_1).T_2}
\quad \textsc{[T-Lam]}
\]

\[
\frac{\Gamma \vdash f : \Pi(x : T_1).T_2
      \quad
      \Gamma \vdash a : T_1}
{\Gamma \vdash f\, a : T_2[x := a]}
\quad \textsc{[T-App]}
\]

\subsection{Open- and Closed-World Variable Rules}

\[
\frac{(x : T) \in \Gamma}
{\Gamma \vdash x : T}
\quad \textsc{[T-Var]}
\]

\[
\frac{x \notin \Gamma \quad \Gamma\ \text{is open-world}}
{\Gamma \vdash x : \mathrm{Unit}}
\quad \textsc{[T-Var-Open]}
\]

\[
\frac{x \notin \Gamma \quad \Gamma\ \text{is closed-world}}
{\Gamma \nvdash x : T \text{ for any } T}
\quad \textsc{[T-Var-Closed]}
\]

\section{Principal Theorems}

\begin{theorem}[Collapse Soundness]
\label{thm:collapse-soundness}
If $\Gamma \vdash \mathrm{collapse}(t, c) : T$ then
$\Gamma \vdash t : \Adm{T'}$ for some $T'$, and $T = \Collapsed{T'}{c}$.
\end{theorem}

\begin{proof}
By inversion on rule \textsc{[T-Collapse]}: the only rule that derives
a Collapse judgement requires the premise $\Gamma \vdash t : \Adm{T}$,
so the inner term must have an admissible type,
and the collapse type is $\Collapsed{T}{c}$ by definition.
\end{proof}

\begin{theorem}[Type Preservation]
\label{thm:preservation}
If $\Gamma \vdash t : T$ and $t \to t'$, then $\Gamma \vdash t' : T$.
\end{theorem}

\begin{proof}[Proof sketch]
By structural induction on the reduction relation $\to$.
The key cases are \textsc{[T-Pop]}, \textsc{[T-Refuse]}, and \textsc{[T-Collapse]},
each of which maps a reduction step to an admissibility transformation
that is closed under the respective rule.
Beta-reduction $(\lambda x.b)\, a \to b[x := a]$ requires a
substitution lemma ($\Gamma, x:T_1 \vdash b:T_2$ and $\Gamma \vdash a:T_1$
implies $\Gamma \vdash b[x:=a]:T_2[x:=a]$) and a canonical forms lemma,
whose proofs we omit; the current implementation realises only the core
fragment and these results serve as proof obligations for future work.
\end{proof}

\begin{theorem}[Progress]
\label{thm:progress}
If $\vdash t : T$ (well-typed in the empty context) then either:
\begin{enumerate}
\item[(i)] $t$ is a value, or
\item[(ii)] there exists $t'$ such that $t \to t'$.
\end{enumerate}
\end{theorem}

\begin{proof}
By structural induction on $t$.
Variable terms are values.
Lambda terms are values (closures).
$\mathrm{pop}$, $\mathrm{refuse}$, $\mathrm{bind}$, and $\mathrm{collapse}$
of a well-typed inner term either reduce (if the inner term reduces)
or produce a value (if the inner term is a value).
\end{proof}

% ─────────────────────────────────────────────────────────────────────────────
\chapter{Open Worlds and Closed Worlds}

\epigraph{\itshape
The type checker's behavior depends not only on the term but on
the assumed boundary of the world.}{---}

\section{The Discovery}

During implementation of the Spherepop type checker, a conflict emerged.
The interpreter treated free variables as symbolic atoms:
an unbound $x$ evaluates to \texttt{Value::Symbol("x")}.
The type checker, in its first version, rejected unbound variables as errors.

Programs that evaluated correctly consistently failed type-checking.
The apparent ``bug fix'' was to make the type checker also treat free variables
as having type \texttt{Unit}.

But this apparent bug conceals a genuine theoretical question:
\emph{which behavior is correct?}
The answer is that both are correct, under different assumptions
about the boundary of the world.

\section{Formal Definitions}

\begin{definition}[Closed-World Context]
A \emph{closed-world context} $\Gamma_c = \{x_1 : T_1, \ldots, x_n : T_n\}$
satisfies:
\[
x \notin \Gamma_c \;\Rightarrow\; \Gamma_c \vdash x : \bot.
\]
The absence of a declaration is the absence of the object.
Rule \textsc{[T-Var-Closed]} applies.
\end{definition}

\begin{definition}[Open-World Context]
An \emph{open-world context} $\Gamma_o = \Gamma_c \cup U$
where $U$ is the \emph{unresolved possibility manifold}
(symbols that have not been declared but have not been refused) satisfies:
\[
x \notin \Gamma_c \;\Rightarrow\; \Gamma_o \vdash x : \mathrm{Unit}.
\]
The absence of a declaration is undecided possibility.
Rule \textsc{[T-Var-Open]} applies.
\end{definition}

\begin{proposition}[Monotonicity of Provability]
\label{prop:ow-monotone}
\[
\Gamma_c \subseteq \Gamma_o
\quad\Rightarrow\quad
\mathrm{Provable}(\Gamma_c) \subseteq \mathrm{Provable}(\Gamma_o).
\]
\end{proposition}

\begin{proof}
Every typing derivation valid under $\Gamma_c$ is valid under $\Gamma_o$,
since $\Gamma_o$ extends $\Gamma_c$ and adds only the open-world variable rule,
which derives additional types for previously-error terms.
\end{proof}

\begin{remark}
The open/closed world distinction is not an engineering convenience.
It is the mechanization of an epistemological commitment:
a closed-world type checker assumes complete knowledge of the reachable;
an open-world type checker assumes partial knowledge.
A language whose fundamental ontology is about possibility space
cannot leave this distinction implicit.
\end{remark}

\begin{implbox}{Implementation consequence: \texttt{types/checker.rs}}
\begin{lstlisting}
pub enum TypeMode {
    ClosedWorld,  // unknown variable is a type error
    OpenWorld,    // unknown variable has type Unit
}
pub struct Context {
    bindings: HashMap<String, Type>,
    pub mode: TypeMode,
}
\end{lstlisting}
The mode is \emph{preserved through \texttt{extend}}: a closed-world context
cannot silently become open-world mid-derivation.
\end{implbox}

% ═══════════════════════════════════════════════════════════════════════════════
\part{Collapse as Quotient}
% ═══════════════════════════════════════════════════════════════════════════════

\chapter{Observable State as a Quotient of History}

\section{Motivation}

The original Spherepop design represented collapse as:
\begin{lstlisting}
Event::Collapse(Symbol)
\end{lstlisting}
A collapse event was simply a record that a symbol had been committed.
This was insufficient, and the insufficiency is philosophically illuminating.

\emph{Observable state depends on how you look.}
The same history, examined by different observers with different
observational frameworks, yields different observable states.
A collapse without a rule is a collapse into an unspecified notion of visibility.

\section{Definitions}

\begin{definition}[Collapse Rule]
\label{def:collapse-rule}
A \emph{collapse rule} is a function $c : \calH \to O_c$
from the history monoid to an \emph{observational space} $O_c$.
\end{definition}

\begin{definition}[Observational Equivalence]
\label{def:obs-equiv}
Two histories $H_1, H_2 \in \calH$ are \emph{observationally equivalent
under rule $c$}, written $H_1 \sim_c H_2$, if and only if
$c(H_1) = c(H_2)$.
\end{definition}

\begin{definition}[Observable State as Quotient]
\label{def:obs-quotient}
The \emph{observable state space} under rule $c$ is the quotient:
\[
O_c = \calH \big/ {\sim_c}.
\]
An observable state is an equivalence class of histories.
\end{definition}

\begin{remark}
Definition~\ref{def:obs-quotient} is the mathematical heart of the monograph.
Observable state is \emph{not} a state.
Observable state is a quotient.
Two histories that are globally distinct may be observationally indistinguishable
under a specific collapse rule.
This is a feature, not a deficiency: the collapse rule specifies which
distinctions are relevant to a given question.
\end{remark}

\section{The Four Canonical Rules}

\begin{definition}[Identity Rule]
$c_I : \calH \to \calH$ is the identity: $c_I(H) = H$.
Under Identity, every event is visible and $O_I \cong \calH$.
\end{definition}

\begin{definition}[LastWrite Rule]
$c_{LW}$ maps $H$ to the most recent pop for each symbol,
with subsequent refusals retroactively removing symbols:
\[
c_{LW}(H) = \{x \mapsto 1 \mid \Pop(x) \in H
  \text{ and } \nexists r. \Rfuse(x,r) \in H \text{ after the last } \Pop(x)\}.
\]
\end{definition}

\begin{definition}[Accumulate Rule]
$c_{\mathrm{acc}}(H)$ maps each symbol to the count of its Pop events:
\[
c_{\mathrm{acc}}(H)(x) = |\{i \mid e_i = \Pop(x)\}|.
\]
\end{definition}

\begin{definition}[Projection Rule]
For a label prefix $L$, the \emph{projection rule} $\pi_L$ restricts to
events whose symbol names share prefix $L$:
\[
\pi_L(H) = [e_i \in H \mid \mathrm{sym}(e_i) \text{ starts with } L].
\]
\end{definition}

\begin{proposition}[Coarsening Order]
The collapse rules form a partial order by coarsening:
$c_1 \leq c_2$ iff $\sim_{c_2} \;\subseteq\; \sim_{c_1}$
(finer equivalence = more informative observation).
Under this order: $c_I$ is the finest (most informative) rule;
$\pi_L$ is coarser than $c_I$; $c_{LW}$ and $c_{\mathrm{acc}}$ are
incomparable in general.
\end{proposition}

\section{Observational Equivalence Theorem}

\begin{theorem}[Observational Equivalence]
\label{thm:obs-equiv}
For any collapse rule $c$ and histories $H_1, H_2$:
$H_1 \sim_c H_2$ if and only if $c(H_1) = c(H_2)$.
Two programs whose interpreter and VM executions yield
$c$-equivalent histories are \emph{observationally indistinguishable}
under rule $c$.
\end{theorem}

\begin{corollary}[History Equivalence Implies Observational Equivalence]
If $I(p) = V(C(p))$ (Replay Equivalence, Theorem~\ref{thm:replay}),
then for every collapse rule $c$:
\[
c(I(p)) = c(V(C(p))).
\]
\end{corollary}

\begin{implbox}{Implementation consequence: \texttt{core/observable.rs}}
\begin{lstlisting}
pub fn derive(history: &History, rule: &CollapseRule)
    -> ObservableState
pub fn equivalent(h1: &History, h2: &History,
                  rule: &CollapseRule) -> bool {
    Self::derive(h1, rule) == Self::derive(h2, rule)
}
\end{lstlisting}
Observable state is \emph{never stored} in the World. It is always
freshly derived. The function \texttt{observe} is a method, not a field,
mechanizing the philosophical claim that observable state is a quotient.
\end{implbox}

% ═══════════════════════════════════════════════════════════════════════════════
\part{Lambda Calculus as Derived Structure}
% ═══════════════════════════════════════════════════════════════════════════════

\chapter{Beta-Reduction as History Contraction}

\section{Motivation}

In the standard account, computation is explained in terms of functions.
Every other computational phenomenon is a deviation from or extension
of the basic function-application model.

Spherepop reverses this priority.
Functions are not primitive.
They are a convenient notation for a particular pattern in history:
\emph{deferred commitment}.

\section{Beta-Reduction as Collapse}

\begin{definition}[Lambda History Encoding]
A lambda abstraction $\lambda x.M$ applied to argument $N$ generates the
history fragment:
\[
[\mathrm{LamIntro}(x),\, \mathrm{Apply}(N)].
\]
The application commits $x$ to $N$, yielding the continued history
beginning with $[\Pop(x)]$ (in the symbolic history of the parameter's use).
\end{definition}

\begin{proposition}[Beta as Quotient]
\label{prop:beta-quotient}
Beta reduction
\[
(\lambda x.M)\,N \;\to\; M[x := N]
\]
is the projection of the history
$[\mathrm{LamIntro}(x), \mathrm{Apply}(N)]$
onto its observable state under the substitution collapse rule
$c_{\mathrm{subst}}$, defined by
$c_{\mathrm{subst}}(H)(x) = N$ whenever $\mathrm{Apply}(N) \in H$
after $\mathrm{LamIntro}(x)$.
\end{proposition}

\begin{remark}
Proposition~\ref{prop:beta-quotient} gives a precise content to the informal
claim that ``beta reduction is a special case of collapse.''
Function application is commitment: $x$ can no longer range over any value
other than $N$.
The application event is a Pop of the parameter $x$ from the option space
of its possible values.
The reduced term $M[x := N]$ is the observable state of the application history
under $c_{\mathrm{subst}}$.
\end{remark}

\section{Closures as Frozen Possibilities}

\begin{definition}[Closure]
A \emph{closure} is a pair $(\lambda x.M, \rho)$ where $\rho$ is the
environment (a partial history of binding events) at the time of lambda introduction.
\end{definition}

\begin{remark}
A closure is a \emph{frozen possibility}: a partial future that will generate
events when applied, suspended in a form that can be passed and stored.
The environment $\rho$ is the reachability context frozen at the moment
of lambda introduction --- the option space and historical commitments
present when the abstraction was formed.
\end{remark}

\begin{implbox}{Implementation consequence: \texttt{eval/value.rs}}
\begin{lstlisting}
Value::Closure {
    param: Symbol,
    param_ty: Type,
    body: Box<Term>,
    env: Env,        // frozen reachability context
}
\end{lstlisting}
The \texttt{env} field is the mechanization of ``frozen possibilities.''
It carries the historical commitments made at closure creation time.
When the closure is applied, this context is restored.
\end{implbox}

% ═══════════════════════════════════════════════════════════════════════════════
\part{Compiling Histories}
% ═══════════════════════════════════════════════════════════════════════════════

\chapter{The Failure of Traditional Compiler Correctness}

\section{Motivation}

The standard notion of compiler correctness is:
\begin{center}
\emph{A compiler is correct if the compiled program produces the same output as
the interpreted program for every input.}
\end{center}
This criterion is adequate for state-centric languages.
It is insufficient for Spherepop.

Two programs that produce the same final value may have constructed
radically different histories to get there,
and those histories are part of what the program \emph{is}.

\section{History Equivalence as Correctness}

\begin{definition}[Interpreter and VM]
Let $I : P \to \calH$ be the interpreter function mapping programs to histories,
and $V : \mathrm{IR} \to \calH$ the VM mapping event-IR blocks to histories,
and $C : P \to \mathrm{IR}$ the compiler.
\end{definition}

\begin{definition}[Spherepop Compiler Correctness]
\label{def:correctness}
A compiler $C$ is \emph{correct} if for all programs $p$ and
initial worlds $W$:
\[
I(p, W).\mathrm{history} = V(C(p), W).\mathrm{history}.
\]
Output equivalence (same final value) is a strictly weaker criterion
that is insufficient for Spherepop.
\end{definition}

\begin{theorem}[Replay Equivalence]
\label{thm:replay}
The Spherepop compiler satisfies Definition~\ref{def:correctness}
for all programs expressible in the Spherepop core calculus
(Refuse, Bind, Collapse, and Seq compositions).
\end{theorem}

\begin{proof}[Proof sketch]
By structural induction on program terms.
Base cases: for each primitive operation (Refuse, Bind, Collapse, Pop),
the compiler emits the corresponding \texttt{IrEvent} variant,
and the VM step for that variant performs exactly the same World mutation
as the interpreter's eval case for the corresponding \texttt{Term} variant.
The history appended is therefore identical.
Inductive case (Seq): each sub-term is compiled in order,
and since histories compose by concatenation (the monoid operation),
the composed history of the compiled form equals the composed history
of the interpreted form by the inductive hypothesis.
\end{proof}

\begin{remark}
Theorem~\ref{thm:replay} is not merely an engineering correctness result.
It is the mechanization of the claim that programs \emph{are} their histories.
The test suite verifies it as a first-class property:
\begin{lstlisting}
assert_eq!(
    world_interp.history,
    world_compiled.history,
);
\end{lstlisting}
Every passing instance of this assertion is a mechanized proof obligation.
\end{remark}

% ─────────────────────────────────────────────────────────────────────────────
\chapter{Event IR: Instructions That Are Not Instructions}

\section{Motivation}

Conventional intermediate representations encode instructions that
mutate machine state.
Spherepop IR encodes \emph{events} --- records of occurrences
that persist as part of the history the VM is constructing.

\begin{definition}[Event IR]
The Spherepop intermediate representation is a sequence of \texttt{IrEvent} terms:
\[
\mathrm{IR} ::= \mathrm{Pop} \;\mid\; \mathrm{Refuse}(r) \;\mid\;
  \mathrm{Collapse}(c) \;\mid\; \mathrm{Bind} \;\mid\;
  \mathrm{Load}(x) \;\mid\; \mathrm{Apply} \;\mid\; \cdots
\]
Each event carries its reason or rule as a first-class argument,
preserving the full certificate structure of the source term.
\end{definition}

\begin{proposition}[IR Faithfulness]
The compilation function $C : \mathrm{Term} \to \mathrm{IR}$ is
\emph{faithful}: for every primitive term $t$,
$C(t)$ contains exactly the event variants that $I(t)$ would append
to the history, with identical arguments.
\end{proposition}

\begin{proof}
By inspection of the compiler's \texttt{lower} function and the
interpreter's \texttt{eval} function:
each branch of \texttt{lower} emits the \texttt{IrEvent} variant that
corresponds bijectively to the \texttt{Event} appended by the matching
\texttt{eval} branch.
\end{proof}

\begin{remark}
The IR Faithfulness Proposition is what makes Theorem~\ref{thm:replay}
(Replay Equivalence) provable. The VM does not ``execute'' IR events;
it \emph{reconstructs} the history that the interpreter would have built.
The VM is a history reconstruction engine, not a state mutation engine.
\end{remark}

% ═══════════════════════════════════════════════════════════════════════════════
\part{Category Theory and Deep Structure}
% ═══════════════════════════════════════════════════════════════════════════════

\chapter{Histories as Morphisms}

\section{The History Category Revisited}

We introduced $\Hist$ in Chapter~3 as the category of option spaces
and histories.
We now develop its structure more carefully in relation to observable state.

\begin{definition}[Option-Space Transition Sequences]
A \emph{transition sequence} from $\Omega$ to $\Omega'$ is a history $H$
such that executing $H$ starting from $\Omega$ yields $\Omega'$.
This is the morphism $H : \Omega \to \Omega'$ in $\Hist$.
\end{definition}

\begin{proposition}[$\Hist$ as a Free Strict Monoidal Category]
$(\Hist, \oplus, \varepsilon, \otimes)$ is a free strict monoidal category where:
\begin{itemize}
\item Tensor product $H_1 \otimes H_2$ is the Meld operation
  (parallel composition of histories).
\item Monoidal unit is the empty option space with empty history.
\item Sequential composition is $\oplus$ (Definition~\ref{def:hist-category}).
\end{itemize}
\end{proposition}

% ─────────────────────────────────────────────────────────────────────────────
\chapter{Collapse Functors}

\section{Collapse as a Functor $\Hist \to \Obs_c$}

\begin{definition}[Observable-State Category $\Obs_c$]
For a collapse rule $c$, define the category $\Obs_c$:
\begin{itemize}
\item \emph{Objects:} observable states $o \in O_c$.
\item \emph{Morphisms:} $\mathrm{Hom}(o_1, o_2)$ is the set of observable
  transitions that transform $o_1$ into $o_2$ under rule $c$.
\item \emph{Composition:} sequential observable transition.
\end{itemize}
\end{definition}

\begin{definition}[Collapse Functor]
\label{def:collapse-functor}
For a collapse rule $c$, define the \emph{collapse functor}
$F_c : \Hist \to \Obs_c$ by:
\[
F_c(\Omega) = c(\varepsilon_\Omega) \quad \text{(object map)}
\]
where $\varepsilon_\Omega$ denotes the empty history at object $\Omega$
(i.e.\ the identity morphism at $\Omega$ in $\Hist$).
\[
F_c(H) = c(H) \quad \text{(morphism map)}
\]
\end{definition}

\begin{theorem}[Functoriality of Collapse]
\label{thm:functor}
$F_c$ is a well-defined functor when the collapse rule $c$
respects composition, i.e., when
\[
c(H_2 \hconcat H_1) = c(H_2) \circ c(H_1)
\]
(up to the appropriate notion of composition in $\Obs_c$).
\end{theorem}

\begin{proof}
Identity preservation: $F_c(\varepsilon) = c(\varepsilon) = \mathrm{id}_{O_c}$
since the empty history yields the empty observable state,
which is the identity transformation in $\Obs_c$.

Composition preservation:
\[
F_c(H_2 \circ H_1) = F_c(H_1 \hconcat H_2) = c(H_1 \hconcat H_2).
\]
When $c$ respects composition:
$c(H_1 \hconcat H_2) = c(H_2) \circ c(H_1) = F_c(H_2) \circ F_c(H_1)$.
\end{proof}

\begin{remark}
Theorem~\ref{thm:functor} elevates collapse from an implementation detail
to a category-theoretic object.
Collapse is not merely ``reading off a value.''
It is a structure-preserving map between categories.
The collapse rule specifies \emph{which structures are preserved} ---
and different rules preserve different structures.
\end{remark}

\begin{remark}
The four canonical rules have different functoriality properties:
$c_I$ (Identity) is trivially functorial (it is the identity functor on $\Hist$).
$c_{\mathrm{acc}}$ (Accumulate) is functorial when the observational composition
is addition of count-maps.
$\pi_L$ (Projection) is functorial: projection of a concatenated history
equals the concatenation of the projected histories.
$c_{LW}$ (LastWrite) is \emph{not} generally functorial, since the last-write
of a concatenated history is not the composition of the last-writes of its parts
--- a later segment can retroactively affect the observable state of an earlier one.
\end{remark}

% ═══════════════════════════════════════════════════════════════════════════════
\part{Representation and Reachability}
% ═══════════════════════════════════════════════════════════════════════════════

\chapter{Filters Create Worlds}
\label{chap:filters}

\epigraph{\itshape
Every representation creates a world with its own notions of sameness,
difference, adjacency, and reachability.}{---}

\section{Motivation}

The Spherepop collapse framework naturally connects to a broader claim about
representations:
\emph{representations do not merely preserve or lose information ---
they reorganize the topology of distinguishability.}

\section{The Santoro–Waghmare–Panaretos Connection}

We interpret a recent result in functional data analysis through the lens
of admissibility geometry.

Let $P, Q \in \mathcal{P}(X)$ be probability measures on a separable metric space.
The kernel covariance embedding
\[
P \;\mapsto\; S_P \;\mapsto\; \mathcal{N}(0, S_P)
\]
maps distributions to Gaussian measures on a Hilbert space.

The result of Santoro, Waghmare, and Panaretos (2024) shows that
for distinct continuous distributions,
$\mathcal{N}(0, S_P)$ and $\mathcal{N}(0, S_Q)$ are mutually singular.
The standard reading is that infinite-dimensional representations amplify
tiny distinctions into maximal separations.

\begin{definition}[Reachability Class]
For a Gaussian measure $\mathcal{N}(0, S_P)$, define its
\emph{reachability class}:
\[
R(P) = \bigl\{Q : \mathcal{N}(0, S_Q) \sim \mathcal{N}(0, S_P)\bigr\}
\]
where $\sim$ denotes measure equivalence (mutual absolute continuity).
\end{definition}

\begin{theorem}[Topological Separation via Representation]
\label{thm:separation}
Under the conditions of the Santoro--Waghmare--Panaretos theorem,
for distinct continuous distributions $P \neq Q$:
\[
d(P, Q) \ll 1
\quad \text{(metric closeness downstairs)}
\]
yet
\[
R(P) \cap R(Q) = \varnothing
\quad \text{(topological separation upstairs).}
\]
The kernel embedding converts metric closeness into topological separation.
\end{theorem}

\begin{remark}
Theorem~\ref{thm:separation} is related to the Feldman--Hájek theorem,
which establishes that two Gaussian measures are either equivalent ($\mu \sim \nu$)
or mutually singular ($\mu \perp \nu$) --- there is no continuum.
In the reachability class language: $R(P) = \{P\}$ for almost all $P$.
\end{remark}

\section{The Representation-Ontology Reading}

The standard interpretation of Theorem~\ref{thm:separation} is epistemological:
\emph{distinct distributions become perfectly distinguishable.}

The CLIO/Spherepop interpretation is ontological:
\emph{representation can reorganize possibility space itself.}

\begin{proposition}[Collapse Rules as Reachability Filters]
Spherepop collapse rules are computationally analogous to the kernel embedding
in Theorem~\ref{thm:separation}.
The rule $c$ maps the history space $\calH$ to an observable state space $O_c$,
changing which histories are distinguishable.
Two histories in the same equivalence class $[H]_{\sim_c}$ are
\emph{topologically identified} under rule $c$:
they occupy the same stratum of observable state space.
\end{proposition}

\begin{remark}
The parallel is precise.
Before applying a collapse rule:
histories may be metrically close (they differ in few events)
yet topologically distinct (they produce different observable states under $c$).
Conversely: histories that are metrically distant (many events differ)
may be topologically identified (they produce the same observable state under $c$).
The collapse rule, like the kernel embedding, changes \emph{adjacency},
not just distance.
\end{remark}

\section{CLIO Projections and History Strata}

\begin{definition}[CLIO Projection]
% CLIO = Constraint-Limited Information Ontology (Flyxion, 2023--).
A CLIO projection $\pi : X \to M$ maps a representation space $X$
to a manifold $M$ of observational strata.
In the Spherepop setting:
$X = \calH$ (history space),
$M = O_c$ (observable state space under rule $c$),
and $\pi = F_c$ (the collapse functor).
\end{definition}

\begin{remark}
Observable states under Projection rule $\pi_L$ are precisely the CLIO strata
corresponding to the label $L$.
Histories become trajectories through stratum space.
The collapse functor $F_c$ is the formalization of CLIO's projection $\pi$.
\end{remark}

% ═══════════════════════════════════════════════════════════════════════════════
\part{Programs as Reachability Sculptures}
% ═══════════════════════════════════════════════════════════════════════════════

\chapter{The Codebase as a Minimal Working Ontology}

\section{Each Module as Ontological Necessity}

The Spherepop codebase is organized to make philosophical dependencies visible.
The following table summarizes the justification for each module:

\medskip
\begin{center}
\begin{tabular}{lp{0.58\textwidth}}
\toprule
Module & Ontological Necessity \\
\midrule
\texttt{core/event.rs}
& Events are primitive. They cannot be derived from types or values
  without circularity. The codebase must begin here. \\[4pt]
\texttt{core/history.rs}
& History is computational identity.
  The free monoid structure (append, meld; no remove) is the mechanization
  of irreversibility. History equality is the primary identity criterion. \\[4pt]
\texttt{core/option\_space.rs}
& Possibility is a runtime structure, not a static specification.
  The conservation law requires both History and OptionSpace to coexist. \\[4pt]
\texttt{core/observable.rs}
& Observable state is derived, not primary.
  Separating it from History enforces the metaphysical priority claim. \\[4pt]
\texttt{types/ty.rs}
& Types must reflect admissibility structure.
  \texttt{Admissible}, \texttt{Refused}, and \texttt{Collapsed}
  are reachability certificates, not value sets. \\[4pt]
\texttt{types/checker.rs}
& Admissibility must be verified.
  \texttt{TypeMode} is the mechanization of the open/closed world epistemology. \\[4pt]
\texttt{eval/interpreter.rs}
& Evaluation is historical construction, not value computation.
  Returning \texttt{EvalResult\{value, steps, admissible\}} is forced. \\[4pt]
\texttt{compiler/lower.rs}
& Compilation must preserve possibility structure.
  Emitting \texttt{IrEvent} rather than \texttt{IrInstruction} is forced. \\[4pt]
\texttt{ir/event\_ir.rs}
& The IR must be event-based because the target runtime is history-based.
  \texttt{Refuse(reason)} and \texttt{Collapse(rule)} carry certificates. \\[4pt]
\texttt{vm/machine.rs}
& Execution is replay, not transformation.
  The VM reconstructs history; correctness is history identity. \\
\bottomrule
\end{tabular}
\end{center}

\section{The Implemented Milestones}

\begin{description}
\item[v0.1 — History-first execution.]
World = History + OptionSpace; evaluation returns EvalResult alongside a World;
the conservation law is enforced as a hard runtime invariant.

\item[v0.2 — Documented refusal.]
\texttt{RefusalReason} is introduced as a first-class type,
threading through events, types, values, and IR.
Refusal is structured inadmissibility with a permanent certificate.

\item[v0.3 — Collapse as quotient.]
\texttt{CollapseRule} specifies the observational framework.
Observable state is derived by \texttt{ObservableState::derive(history, rule)}.
Observational equivalence becomes rule-relative and formally testable.

\item[v0.4 — TypeMode.]
\texttt{TypeMode::OpenWorld} and \texttt{TypeMode::ClosedWorld} are introduced.
The type checker's behavior depends on this mode,
mechanizing the philosophical discovery about world boundaries.

\item[v0.5 — History equivalence as compiler correctness.]
The correctness criterion $I(p).\mathrm{history} = V(C(p)).\mathrm{history}$
is implemented as a first-class test obligation.
\end{description}

% ─────────────────────────────────────────────────────────────────────────────
\chapter{The Thesis Defended}

\epigraph{\itshape
Computation is the progressive restriction of possibility,
not the manipulation of values.}{---}

\section{The Outrageous Claim Made Precise}

The thesis sounds outrageous only as a slogan.
When a programmer writes \texttt{x = 5},
they are not merely assigning a value.
They are committing to a specific future ---
one in which the value of \texttt{x} is 5 ---
and foreclosing all futures in which $x$ has any other value.
The assignment is a Pop: a possibility removed from the option space.

When they write an \texttt{if}-\texttt{else} branch,
they are partitioning possibility space:
the condition is the criterion; the branches are the strata.

When they write a function call,
they are deferring a commitment:
the function is a frozen possibility,
and the call is the moment of actualization.

The Spherepop framework makes this geometry \emph{explicit}.
Pop, Refuse, Bind, and Collapse are the explicit operations on possibility space
that conventional languages hide behind assignment, branching, and function calls.

\begin{proposition}[Expressive Correspondence]
The following conventional constructs are expressible in terms of
Spherepop primitives:
\begin{itemize}
\item \emph{Variable assignment} $x := v$ $\approx$ $\mathrm{Pop}(v)$ followed by $\mathrm{Bind}(x, v)$.
\item \emph{Conditional branch} on predicate $\phi$ $\approx$ $\mathrm{Pop}(\phi^+)$ or $\mathrm{Refuse}(\phi^+, \mathrm{ConstraintViolation}(\phi))$ followed by $\mathrm{Pop}(\phi^-)$.
\item \emph{Function application} $f(a)$ $\approx$ $\mathrm{LamIntro}(x)$ followed by $\mathrm{Apply}(a)$ (a deferred Pop).
\item \emph{Exception throwing} $\approx$ $\mathrm{Refuse}(t, r)$ (documented inadmissibility without system destruction).
\end{itemize}
\end{proposition}

\section{Why This Changes Everything}

In the conventional account, programs are descriptions of computations:
sequences of instructions that tell the machine what to do.

In the Spherepop account, programs are \emph{constructions within possibility space}:
paths through the geometry of accessible futures.

The difference is not cosmetic.
It affects what questions we can ask:

\begin{center}
\begin{tabular}{ll}
\toprule
Conventional question & Spherepop question \\
\midrule
What is the current state? & What remains reachable? \\
What is the output value? & What history was constructed? \\
Did the program terminate? & What possibilities were foreclosed? \\
Is there a bug? & Which commitment was inadmissible? \\
What did the function return? & Under what rule is the result observable? \\
\bottomrule
\end{tabular}
\end{center}

The Spherepop questions are not more abstract.
They are more \emph{honest}: they expose what computation actually does
to the geometry of possibility, rather than summarizing the result
as a degenerate projection.

% ─────────────────────────────────────────────────────────────────────────────
\chapter{Future Directions}
\label{chap:future-dir}

\section{Proof-Carrying Refusal}

The \texttt{RefusalReason} type is a vocabulary for inadmissibility.
It is not yet a proof system.
A fully realized Spherepop type system would require \texttt{RefusalReason}
to be a proof term: evidence that can be checked, not merely a label.
This connects to the theory of certified refusal:
a refusal is well-formed only if accompanied by a checkable argument
for why the refused path was inadmissible.

\section{Dependent Process Types}

The current $\mathrm{Process}(T_1 \rightsquigarrow T_2)$ type is first-order.
A dependent process type
$\Pi(x : T_1).\,\mathrm{Process}(x, f(x))$
would allow the type of the output to depend on the specific value consumed.
This is the natural extension of dependent type theory into the process calculus setting.

\section{Admissibility Lattices}

The current admissibility structure is Boolean: a term is admissible or not.
A more refined theory would recognize that admissibility is graded ---
some futures are more constrained than others,
and the constraints form a lattice.
This connects to the broader admissibility geometry program of RSVP and CLIO.

\section{Collapse Functor Composition}

The four canonical collapse rules are special cases of a more general
family of functors $F_c : \Hist \to \Obs_c$.
A theory of \emph{functor composition} would characterize which combinations
of rules are valid --- for instance, when
$F_{c_2} \circ F_{c_1}$ is itself a valid collapse functor.

\section{Distributed Spherepop Runtimes}

The current World is single-threaded.
A distributed Spherepop runtime would coordinate multiple Worlds,
each with their own histories, under a global consistency constraint.
The consistency constraint must be expressed in terms of history equivalence
under a shared collapse rule ---
not value consistency.
The event sourcing and CRDT literatures provide partial models;
the Spherepop-specific challenge is the certified-refusal structure.

% ─────────────────────────────────────────────────────────────────────────────

\chapter*{Coda: The Implementation as Philosophical Argument}
\addcontentsline{toc}{chapter}{Coda: Implementation as Argument}

\epigraph{\itshape
The implementation did not merely realize the theory;
it exposed which distinctions the theory was still hiding.}{---}

The central claim of this monograph is not that Spherepop is a good
programming language, though we believe it is an interesting one.
The central claim is that \emph{taking a philosophical position seriously
enough to implement it is a form of philosophical argument.}

The Spherepop implementation is not a demonstration of a pre-formed theory.
It is a test of a theory's coherence.
The theory said that histories are primary, types are admissibility certificates,
and refusal is documented inadmissibility.
The implementation asked: what does this actually mean in executable form?

The four implementation discoveries were not predicted by the initial framework.
They were forced by the attempt to mechanize it.

The original \texttt{Refuse(Symbol)} gestured at documented inadmissibility.
It took implementation pressure to reveal that a symbol alone documents nothing:
the reason is the document.

The original notion of observable state gestured at the quotient structure.
It took \texttt{World::observe} to force the question: observable under what rule?

The original type checker gestured at world-relative knowledge.
It took the free-variable conflict to force the question:
which world are we assuming?

The original correctness criterion gestured at semantic preservation.
It took the compiler test to force the question:
preservation of what, exactly?

Each discovery reveals a distinction the philosophical framework was
gesturing at without making precise.
Only in the attempt to mechanize --- to force the philosophy to compile ---
did the gestures become definitions.

This is the deepest lesson of ontology-driven language design.
Philosophy can identify the right questions.
Only implementation can force the right precision.

\medskip\noindent
The Spherepop codebase is a philosophical argument written in Rust.
Like all arguments, it became more precise in the act of being made.

% ═══════════════════════════════════════════════════════════════════════════════
% NEW CHAPTERS — inserted between Part VIII and \backmatter
% ═══════════════════════════════════════════════════════════════════════════════

% ─────────────────────────────────────────────────────────────────────────────
\part{Limits, Gaps, and Open Problems}
% ─────────────────────────────────────────────────────────────────────────────

\chapter{The Observation Limit}

\epigraph{\itshape
Spherepop does not abolish the state illusion.
It makes the quotient explicit, parameterised, and reversible
whenever the underlying history is retained.}{---}

\section{Motivation: A Reviewer's Challenge}

A careful reader of this monograph will notice a tension.
The document claims to solve the state illusion ---
the loss of historical information caused by projecting history to state.
Yet every output of a Spherepop program is a \emph{collapse}:
a value derived by applying a rule to a history.
The programmer observes the collapsed value, not the history.
Is Spherepop not simply relocating the state illusion rather than resolving it?

This chapter addresses that challenge directly.
We prove a fundamental limit --- the \emph{No Direct Observation Theorem} ---
that shows the challenge is not a defect in Spherepop but a mathematical
necessity.
We then distinguish precisely between two phenomena the challenge conflates:
\emph{state erasure} and \emph{state illusion}.
Spherepop eliminates the first.
No computational system can eliminate the second.

\section{Definitions}

\begin{definition}[Observation Map]
An \emph{observation map} for a collapse rule $c$ is any function
\[
\phi : \calH \to O
\]
from histories to an observable space $O$.
\end{definition}

\begin{definition}[History-Faithful Observation]
An observation map $\phi$ is \emph{history-faithful} if it is injective:
$\phi(H_1) = \phi(H_2) \Rightarrow H_1 = H_2$.
\end{definition}

\section{The No Direct Observation Theorem}

\begin{theorem}[No Direct Observation]
\label{thm:no-direct-obs}
Let $O$ be any set with $|O| < |\calH|$.
Then every observation map $\phi : \calH \to O$ is non-injective:
there exist distinct $H_1 \neq H_2$ with $\phi(H_1) = \phi(H_2)$.
Consequently, no computational system with a finite observable space
can directly verify the history it constructed.
\end{theorem}

\begin{proof}
Since $\calH = \bigcup_{n \geq 0} E^n$ is countably infinite
and $|O| < |\calH|$,
by the pigeonhole principle $\phi$ cannot be injective.
Therefore there exist $H_1 \neq H_2$ with $\phi(H_1) = \phi(H_2)$.
\end{proof}

\begin{remark}
Theorem~\ref{thm:no-direct-obs} is the fundamental limit of computational observation.
It says: \emph{every observation is a quotient}.
This is not a contingent limitation that better engineering can overcome.
It is a cardinality constraint.
A system that made its observation map injective --- that recorded every
distinction in full history space --- would itself be a history.
To observe computation from outside, one must compress.
Every compression identifies distinct histories.
\end{remark}

\section{Erasure vs Illusion}

The Theorem~\ref{thm:no-direct-obs} suggests a distinction the literature elides:

\begin{definition}[State Erasure]
A computational system suffers \emph{state erasure} if it irrecoverably
discards the history $H$ after computing $\sigma(H)$.
That is, after execution, $H$ is unavailable and only $\sigma(H)$ remains.
\end{definition}

\begin{definition}[State Illusion]
A computational system presents a \emph{state illusion} to its programmers
if they can only observe $\sigma(H)$ and not $H$ directly.
By Theorem~\ref{thm:no-direct-obs}, \emph{some} state illusion is unavoidable
for any non-trivial observable space.
\end{definition}

\begin{proposition}[What Spherepop Eliminates]
\label{prop:what-sp-eliminates}
Spherepop eliminates state erasure:
the history $H$ is retained in \texttt{world.history}
and is available to any observer.

Spherepop cannot and does not eliminate the state illusion:
the programmer's primary observable is $c(H)$, not $H$,
so distinct histories that agree under $c$ remain indistinguishable
from the perspective of that observation.
\end{proposition}

\begin{remark}
The correct claim is therefore:

\begin{center}
\emph{Spherepop does not abolish the state illusion.
It makes the quotient explicit, parameterised, and history-preserving.}
\end{center}

Concretely: two histories $H_1 \sim_c H_2$ are observationally equivalent
under rule $c$, but both are retained.
A programmer who suspects their collapsed observation is misleading
can switch to a finer collapse rule $c'$ (where $\sim_{c'} \subsetneq \sim_c$),
recover additional distinctions, and potentially identify which of many
$c$-equivalent histories they actually constructed.
In conventional systems, this recovery is impossible because $H$ was discarded.
\end{remark}

\section{Rule Refinement as Disambiguation}

\begin{definition}[Rule Refinement Order]
Collapse rule $c_1$ is \emph{finer than} $c_2$, written $c_1 \preceq c_2$, if
$H_1 \sim_{c_1} H_2 \Rightarrow H_1 \sim_{c_2} H_2$:
every distinction preserved by $c_1$ is also preserved by $c_2$.
The Identity rule $c_I$ is the finest (maximally informative) rule;
the trivial rule $c_\top$ mapping all histories to a single point is the coarsest.
\end{definition}

\begin{proposition}[Disambiguation by Refinement]
\label{prop:disambiguation}
Let $H_1 \sim_{c} H_2$ (the histories are indistinguishable under rule $c$)
and $H_1 \neq H_2$.
Then there exists a finer rule $c' \preceq c$ such that
$c'(H_1) \neq c'(H_2)$.
In the limit, $c_I$ always disambiguates since $c_I(H) = H$ exactly.
\end{proposition}

\begin{proof}
Take $c' = c_I$: since $H_1 \neq H_2$ and $c_I(H) = H$,
we have $c_I(H_1) = H_1 \neq H_2 = c_I(H_2)$.
Therefore $c_I$ is a finer rule that disambiguates $H_1$ and $H_2$.
\end{proof}

This proposition formalises the engineering value of history preservation.
A system that erases histories cannot apply a finer rule later.
A Spherepop system that preserves histories can always refine its observation
by switching to a finer collapse rule.

% ─────────────────────────────────────────────────────────────────────────────
\chapter{When Does Observational Equivalence Imply History Equivalence?}

\section{Motivation}

The deepest open question in the Spherepop framework is the converse
of the collapse functor:

\begin{center}
$H_1 = H_2 \Rightarrow c(H_1) = c(H_2)$ for every $c$. \hfill (trivially true)\\[6pt]
$c(H_1) = c(H_2) \Rightarrow H_1 = H_2$? \hfill (when does this hold?)
\end{center}

When does observational equivalence under a collapse rule imply history identity?
Understanding the boundary is the deepest open mathematical problem in the framework
and connects to information theory, inverse problems, and kernel geometry.

\section{Definitions}

\begin{definition}[Observational Completeness]
A collapse rule $c$ is \emph{observationally complete} if it is injective:
$c(H_1) = c(H_2) \Rightarrow H_1 = H_2$.
\end{definition}

\begin{definition}[Collapse Kernel]
The \emph{kernel} of a collapse rule $c$ is
\[
\ker(c) = \{(H_1, H_2) \in \calH \times \calH \mid c(H_1) = c(H_2)\}.
\]
Observational completeness holds iff $\ker(c) = \{(H,H) \mid H \in \calH\}$
(the diagonal).
\end{definition}

\section{Completeness Results for Canonical Rules}

\begin{proposition}[Identity Rule is Observationally Complete]
The Identity rule $c_I$ with $c_I(H) = H$ is observationally complete.
\end{proposition}

\begin{proof}
$c_I(H_1) = H_1$ and $c_I(H_2) = H_2$; if $c_I(H_1) = c_I(H_2)$ then $H_1 = H_2$.
\end{proof}

\begin{proposition}[LastWrite is Not Observationally Complete]
The LastWrite rule $c_{LW}$ is not observationally complete.
\end{proposition}

\begin{proof}
Consider the event alphabet $E = \{\mathtt{pop}(x), \mathtt{pop}(y)\}$.
Let $H_1 = [\mathtt{pop}(x), \mathtt{pop}(y)]$ and
$H_2 = [\mathtt{pop}(y), \mathtt{pop}(x), \mathtt{pop}(y)]$.
Under $c_{LW}$: $c_{LW}(H_1)$ maps $x \mapsto 1, y \mapsto 1$ and
$c_{LW}(H_2)$ maps $x \mapsto 1, y \mapsto 1$ (last pop of $y$ wins).
Thus $c_{LW}(H_1) = c_{LW}(H_2)$ but $H_1 \neq H_2$.
\end{proof}

\begin{proposition}[Accumulate is Not Observationally Complete]
The Accumulate rule $c_{\mathrm{acc}}$ is not observationally complete.
\end{proposition}

\begin{proof}
Let $H_1 = [\mathtt{pop}(x), \mathtt{pop}(y)]$ and
$H_2 = [\mathtt{pop}(y), \mathtt{pop}(x)]$.
Both map to $\{x \mapsto 1, y \mapsto 1\}$ under $c_{\mathrm{acc}}$,
but $H_1 \neq H_2$ (different event ordering).
\end{proof}

\begin{remark}
The Accumulate rule identifies histories that are permutations of one another:
$c_{\mathrm{acc}}(H_1) = c_{\mathrm{acc}}(H_2)$ iff $H_1$ and $H_2$ have
the same multiset of events.
This corresponds to an \emph{order-forgetful} observation ---
one that records what happened but not when, relative to other events.
\end{remark}

\section{The Distinguishability Topology}

\begin{definition}[Collapse-Induced Topology]
For a collapse rule $c : \calH \to O_c$, define the
\emph{$c$-induced topology} on $\calH$ as the initial topology
with respect to $c$:
\[
\tau_c = \{ c^{-1}(U) \mid U \subseteq O_c \}.
\]
This makes $c$ continuous (in the discrete topology on $O_c$)
and $\tau_c$ is the coarsest topology that $c$ is continuous with respect to.
\end{definition}

\begin{definition}[History Distance Under Rule $c$]
\[
d_c(H_1, H_2) = \mathbf{1}[c(H_1) \neq c(H_2)]
\]
is the \emph{observational distance} between $H_1$ and $H_2$ under rule $c$.
This is a pseudometric: $d_c(H_1, H_2) = 0$ iff $H_1 \sim_c H_2$.
\end{definition}

\begin{proposition}[Rule Change Induces Topology Change]
\label{prop:topology-change}
For rules $c_1 \preceq c_2$ (where $c_1$ is finer),
\[
\tau_{c_2} \subseteq \tau_{c_1}.
\]
Refining the collapse rule strictly expands the induced topology on $\calH$,
making more distinctions between histories visible.
\end{proposition}

\begin{proof}
If $c_1 \preceq c_2$, then $H_1 \sim_{c_1} H_2 \Rightarrow H_1 \sim_{c_2} H_2$,
so the equivalence classes of $\sim_{c_2}$ are coarser than those of $\sim_{c_1}$.
Every open set in $\tau_{c_2}$ (a union of $\sim_{c_2}$-classes) is also a union
of $\sim_{c_1}$-classes, hence in $\tau_{c_1}$.
\end{proof}

\begin{remark}
Proposition~\ref{prop:topology-change} is the formal content of the claim
that ``changing the collapse rule changes the topology of distinguishability.''
It connects the Spherepop collapse framework to the kernel embedding discussion
of Chapter~\ref{chap:filters}: in both cases, a change of representation
(kernel embedding; collapse rule) changes the topology on the underlying space
(probability distributions; histories).

The kernel embedding moves pairs $(P, Q)$ from the same topology class
(metric closeness) to different topology classes (mutual singularity).
Collapse rule refinement moves pairs $(H_1, H_2)$ from the same equivalence class
to different equivalence classes.
Both are instances of the same structural operation: representation induces topology,
and changing representation changes which distinctions are topologically visible.
\end{remark}

\section{Open Problem: The Inverse Problem}

\begin{notation}
We write $[H]_c = \{H' \in \calH \mid H' \sim_c H\}$ for the equivalence class
of $H$ under rule $c$.
\end{notation}

The inverse problem for a collapse rule $c$ is:
\emph{given an observed state $o \in O_c$, characterise $c^{-1}(o)$.}

This is the formal version of the question ``which histories could have
produced this observation?''

\begin{definition}[Inverse-Problem Fibre]
For $o \in O_c$, the \emph{fibre} of $o$ is $c^{-1}(o) = \{H \in \calH \mid c(H) = o\}$.
\end{definition}

For the Accumulate rule, $c_{\mathrm{acc}}^{-1}(m)$ is the set of all histories
whose event multiset is $m$ --- an infinite set (any permutation of $m$ extended
by arbitrarily many Refuse or Bind events that don't change counts).

Characterising fibres for general rules, and determining when a fibre contains
a unique history (i.e.\ when the rule is injective on a subclass of histories),
is the primary open mathematical problem in the framework.

% ─────────────────────────────────────────────────────────────────────────────
\chapter{Category Theory as Engine, Not Commentary}
\label{chap:cat-engine}

\section{Motivation}

A reviewer of this work identified that the category theory in
Chapters~11--12 is \emph{descriptive} rather than \emph{generative}:
no theorem currently requires categorical machinery to state or prove.
This chapter corrects that deficiency by establishing results that
are naturally and essentially categorical.

\section{The Collapse Composition Theorem}

The collapse functor framework becomes non-trivial when we ask about
\emph{composition of collapse rules}: what happens when we apply one
collapse to the result of another?

\begin{definition}[Rule Composition]
For rules $c_1 : \calH \to O_1$ and $c_2 : O_1 \to O_2$,
the \emph{composed rule} is
\[
(c_2 \circ c_1)(H) = c_2(c_1(H)).
\]
\end{definition}

\begin{theorem}[Collapse Composition]
\label{thm:collapse-composition}
If $F_{c_1} : \Hist \to \Obs_{c_1}$ and
$F_{c_2} : \Obs_{c_1} \to \Obs_{c_2}$
are both valid collapse functors, then their composition
$F_{c_2} \circ F_{c_1}$ is a valid collapse functor
$F_{c_2 \circ c_1} : \Hist \to \Obs_{c_2 \circ c_1}$,
and
\[
F_{c_2 \circ c_1}(H) = F_{c_2}(F_{c_1}(H)).
\]
\end{theorem}

\begin{proof}
Composition of functors is a functor: identity preservation and composition
preservation follow directly from those properties of $F_{c_1}$ and $F_{c_2}$.
The equation $F_{c_2 \circ c_1}(H) = F_{c_2}(F_{c_1}(H))$ holds by
definition of rule composition and the morphism map of $F_c$.
\end{proof}

\begin{remark}
Theorem~\ref{thm:collapse-composition} establishes collapse rules as a
composable algebraic structure.
A sequence of observations $c_1, c_2, \ldots, c_n$ can be composed
into a single rule $c_n \circ \cdots \circ c_1$.
This is the formal basis for \emph{observational pipelines}: a program
may project history through a sequence of finer-to-coarser rules,
each stripping away distinctions that are irrelevant to the next stage.
\end{remark}

\section{The Observation Adjunction}

The deepest categorical result in the framework is an adjunction
between the functor that \emph{constructs} histories and the functor
that \emph{observes} them.

\begin{definition}[History Construction Functor]
Let $\mathbf{Terms}$ be the category whose objects are Spherepop terms
and whose morphisms are reduction sequences.
Define $G : \mathbf{Terms} \to \Hist$ by
\[
G(t) = I(t).\text{history}
\]
(the history produced by interpreting $t$).
\end{definition}

\begin{definition}[Observation Functor]
For a fixed collapse rule $c$, define $\mathcal{O}_c : \Hist \to \Obs_c$ by
\[
\mathcal{O}_c(H) = c(H) = F_c(H).
\]
\end{definition}

\begin{theorem}[History-Observation Adjunction]
\label{thm:adjunction}
For the Identity rule $c_I$,
$\mathcal{O}_{c_I} \dashv G^{\mathrm{op}}$
in the sense that there is a natural bijection
\[
\Hist(G(t), H) \;\cong\; \mathbf{Terms}(t,\, G^{-1}(H))
\]
whenever $G^{-1}(H)$ is defined (i.e.\ when $H$ is the history of some
term under the interpreter).
\end{theorem}

\begin{proof}[Proof sketch]
Under $c_I$, observation is the identity functor on $\Hist$.
The adjunction then reduces to the statement that morphisms in $\Hist$
from $G(t)$ to $H$ correspond to reduction sequences from $t$ to a term
whose interpretation produces $H$.
This is the categorical statement of the Replay Equivalence
(Theorem~\ref{thm:replay}): the interpreter and compiler produce
histories that are morphisms in $\Hist$.
A complete proof requires a full formalisation of $\mathbf{Terms}$
beyond the scope of the current implementation.
\end{proof}

\section{What Is Now Category-Theoretically Forced}

\begin{proposition}[Theorems Requiring Categorical Language]
The following results are most naturally stated in categorical terms:
\begin{enumerate}
\item Theorem~\ref{thm:collapse-composition}
  (Collapse Composition) --- requires functor composition.
\item Theorem~\ref{thm:adjunction}
  (History-Observation Adjunction) --- requires adjunction.
\item Proposition~\ref{prop:topology-change}
  (Rule Change Induces Topology Change) ---
  requires the category of topological spaces and continuous maps.
\item Theorem~\ref{thm:functor}
  (Functoriality of Collapse) ---
  requires the definition of $\Hist$ as a category.
\end{enumerate}
The category theory is now \emph{generative}:
these theorems cannot be stated, let alone proved,
without the categorical framework introduced in Chapter~\ref{chap:cat-engine}.
\end{proposition}

% ─────────────────────────────────────────────────────────────────────────────
\chapter{Strengthening the Kernel Connection}
\label{chap:kernel-strengthened}

\section{Motivation}

Chapter~\ref{chap:filters} (Filters Create Worlds) presented the connection
between Spherepop collapse rules and the Santoro--Waghmare--Panaretos theorem
as an analogy.
The reviewer correctly identified that the analogy was loose:
the kernel theorem is a \emph{limitative} result (separation is unavoidable);
the collapse framework is a \emph{design freedom} (the programmer chooses the rule).

This chapter sharpens the connection by making it a theorem about
\emph{representation-induced topological change} that applies uniformly
to both settings.

\section{A Unified Framework}

\begin{definition}[Representation Map]
A \emph{representation map} is any function $\rho : X \to Y$ from a source
space $X$ to a target space $Y$, equipped with the initial topology
$\tau_\rho = \{\rho^{-1}(U) \mid U \subseteq Y\}$ on $X$.
\end{definition}

\begin{definition}[Representation-Induced Topology]
For a representation map $\rho : X \to Y$:
\begin{itemize}
\item $x_1 \sim_\rho x_2$ iff $\rho(x_1) = \rho(x_2)$
  (equivalence induced by $\rho$).
\item The \emph{separation power} of $\rho$ is
  $|\{[x]_\rho \mid x \in X\}| = |Y|$
  (the number of equivalence classes, bounded by the size of $Y$).
\end{itemize}
\end{definition}

\begin{theorem}[Representation Changes Topology]
\label{thm:rep-topology}
Let $\rho_1, \rho_2 : X \to Y$ be two representation maps with
$\tau_{\rho_1} \neq \tau_{\rho_2}$.
Then there exist $x_1, x_2 \in X$ such that:
\[
x_1 \sim_{\rho_1} x_2 \quad \text{but} \quad x_1 \not\sim_{\rho_2} x_2
\]
or vice versa.
That is, the two representations disagree on whether $x_1$ and $x_2$
are the same object.
\end{theorem}

\begin{proof}
If $\tau_{\rho_1} \neq \tau_{\rho_2}$, there exists an open set
$U \in \tau_{\rho_1} \setminus \tau_{\rho_2}$ (WLOG).
Then $U = \rho_1^{-1}(V)$ for some $V \subseteq Y$, but
$U \neq \rho_2^{-1}(W)$ for any $W$.
Therefore there exist $x_1, x_2$ with $\rho_1(x_1) = \rho_1(x_2)$
(both in the same fibre of $\rho_1$) but $\rho_2(x_1) \neq \rho_2(x_2)$.
\end{proof}

\begin{remark}
Theorem~\ref{thm:rep-topology} is the general form of the observation that
``representations reorganise topology.''
It applies to:
\begin{itemize}
\item \textbf{Kernel embeddings}:
  $\rho_1 = \mathrm{id}$ (identity on distributions),
  $\rho_2 = \mathcal{N}(0, S_\cdot)$ (Gaussian embedding).
  By Theorem~\ref{thm:separation}, $\tau_{\rho_1}$ and $\tau_{\rho_2}$
  differ: $\rho_2$ separates distributions that $\rho_1$ places in
  the same metric neighbourhood.
\item \textbf{Collapse rules}:
  $\rho_1 = c_{\mathrm{acc}}$, $\rho_2 = c_I$.
  By Proposition~\ref{prop:topology-change},
  $\tau_{c_I} \supsetneq \tau_{c_{\mathrm{acc}}}$:
  the identity rule separates histories that the accumulate rule identifies.
\end{itemize}
Both are instances of Theorem~\ref{thm:rep-topology}.
The connection is now a theorem, not an analogy.
\end{remark}

\begin{remark}
The difference between the kernel case and the collapse case is now also precise:
\begin{itemize}
\item In the kernel case, the topological change is \emph{forced}:
  the Gaussian embedding is the unique representation with the mutual-singularity
  property (by Feldman--Hájek).
  The programmer does not choose whether the separation occurs.
\item In the collapse case, the topological change is \emph{chosen}:
  the programmer selects the collapse rule $c$, thereby choosing
  which distinctions are topologically visible.
\end{itemize}
One is a mathematical limit; the other is a design parameter.
Both are instances of Theorem~\ref{thm:rep-topology},
applied in different regimes of freedom.
\end{remark}

% ─────────────────────────────────────────────────────────────────────────────
% APPENDIX
% ─────────────────────────────────────────────────────────────────────────────
% ═══════════════════════════════════════════════════════════════════════════════
% V3 ADDITIONS — inserted at end of new-chapters, before existing \appendix
% Contains:
%   - Strengthened theorems injected into existing chapters (via \input patch)
%   - Appendix B: Category-Theoretic Foundations
%   - Appendix C: Sheaf-Theoretic Semantics
%   - Appendix D: BNF Grammar
% ═══════════════════════════════════════════════════════════════════════════════

% ─────────────────────────────────────────────────────────────────────────────
\part{Deeper Mathematical Structure}
% ─────────────────────────────────────────────────────────────────────────────

\chapter{Derivations: Structures That Emerge Inevitably}
\label{chap:derivations}

\epigraph{\itshape
A framework matures not when it accumulates more propositions
but when it shows that many seemingly separate propositions
are consequences of the same underlying object.}{---}

\section{Motivation}

The preceding chapters have established definitions and verified
properties one by one.
This chapter takes a different posture: it asks which results
can be \emph{derived} from a small number of primitives,
so that the reader sees not merely that something is true
but that it \emph{had} to be true given the ontological commitments
of the framework.

We present six derivations, each revealing a theorem as
the inevitable consequence of a prior commitment.

\section{Derivation 1: The Generalised Possibility Functional}

Theorem~2.6 (Conservation of Possibility) established
$|H_t| + |\Omega_t| = |\Omega_0|$ under pure Pop dynamics.
That result treats all events equally.
The following derivation reveals the original theorem as a
special case of a weighted conservation measure.

\begin{definition}[Event Weight Function]
Define $w : E \to \mathbb{N}$ by
\[
w(\Pop(x)) = 1, \quad
w(\Rfuse(x,r)) = 0, \quad
w(\Bind(a,b)) = 0, \quad
w(\Col(x,c)) = 0.
\]
\end{definition}

\begin{definition}[Generalised Possibility Functional]
\label{def:gpf}
For a world $(H, \Omega)$, define
\[
\Pi(H, \Omega) = |\Omega| + \sum_{e \in H} w(e).
\]
\end{definition}

\begin{theorem}[Conservation of the Possibility Functional]
\label{thm:gpf-conservation}
For any legal execution sequence
$(H_0, \Omega_0) \to (H_1, \Omega_1) \to \cdots \to (H_n, \Omega_n)$,
\[
\Pi(H_t, \Omega_t) = |\Omega_0|
\]
for all $t \geq 0$.
\end{theorem}

\begin{proof}
At $t=0$: $\Pi(H_0, \Omega_0) = |\Omega_0| + 0 = |\Omega_0|$.
At each step, exactly one of four operators is applied:
\begin{itemize}
\item \textbf{Pop}: $|\Omega|$ decreases by 1; $w(\Pop) = 1$ is added.
  Net change: $(-1) + 1 = 0$.
\item \textbf{Refuse}: $|\Omega|$ unchanged; $w(\Rfuse) = 0$ is added.
  Net change: $0$.
\item \textbf{Bind}: $|\Omega|$ unchanged; $w(\Bind) = 0$ is added.
  Net change: $0$.
\item \textbf{Collapse}: $|\Omega|$ unchanged; $w(\Col) = 0$ is added.
  Net change: $0$.
\end{itemize}
In every case $\Pi(H_{t+1}, \Omega_{t+1}) = \Pi(H_t, \Omega_t)$.
By induction, $\Pi(H_t, \Omega_t) = \Pi(H_0, \Omega_0) = |\Omega_0|$.
\end{proof}

\begin{remark}
Theorem~2.6 follows immediately: setting all weights to 1 and
restricting to Pop-only executions recovers $|H_t| + |\Omega_t| = |\Omega_0|$.
But Theorem~\ref{thm:gpf-conservation} is more general: the functional
$\Pi$ is conserved even when Refuse, Bind, and Collapse are intermixed.
The weight function $w$ makes explicit what counts as ``consuming
possibility'' (Pop, weight 1) versus ``documenting structure without
consuming it'' (Refuse, Bind, Collapse, weight 0).
\end{remark}

\begin{corollary}[Irreversibility of Execution]
\label{cor:irrev}
For any non-empty execution sequence,
$|H_{t+1}| > |H_t|$.
Therefore no legal execution can return to a previous world state
$(H_t, \Omega_t)$.
\end{corollary}

\begin{proof}
Every operator appends exactly one event to $H$ (Definition~2.3),
so $|H_{t+1}| = |H_t| + 1 > |H_t|$.
Since $H$ grows strictly and executions are deterministic,
$(H_{t+1}, \Omega_{t+1}) \neq (H_t, \Omega_t)$.
\end{proof}

\begin{remark}
Corollary~\ref{cor:irrev} establishes that Spherepop worlds form a
\emph{directed acyclic execution graph}: world states are nodes, operators
are edges, and there are no cycles.
This gives a rigorous foundation for the informal claim that Spherepop
histories are intrinsically irreversible.
\end{remark}

\section{Derivation 2: Universal Property of Collapse}

The quotient construction $O_c = \calH / {\sim_c}$ is not merely
a convenient definition.
It has a \emph{universal property} that explains why collapse is the
canonical projection — not one representation among many but the
coarsest representation preserving exactly the distinctions
specified by rule $c$.

\begin{theorem}[Universal Property of Collapse]
\label{thm:universal-collapse}
Let $q_c : \calH \to O_c$ be the quotient map for rule $c$.
If $f : \calH \to X$ is any function satisfying
\[
H_1 \sim_c H_2 \implies f(H_1) = f(H_2)
\]
(i.e.\ $f$ is constant on $c$-equivalence classes),
then there exists a unique $\bar f : O_c \to X$ such that
$f = \bar f \circ q_c$.
\end{theorem}

\begin{proof}
By the universal property of quotient sets.
Define $\bar f([H]_c) = f(H)$.
This is well-defined because $f$ is constant on equivalence classes.
Uniqueness: any $\bar f$ satisfying $f = \bar f \circ q_c$ must map
$[H]_c$ to $f(H)$, so $\bar f$ is uniquely determined.
\end{proof}

\begin{remark}
Theorem~\ref{thm:universal-collapse} is philosophically important.
It says that every observation $f$ that ``respects'' a collapse rule $c$
factors uniquely through $O_c$.
Observable state is therefore not merely one representation among many.
It is the \emph{universal} representation among all representations
that identify the same histories as $c$ does.
Any coarser observation necessarily passes through $O_c$.
\end{remark}

\section{Derivation 3: Unique History Factorisation}

\begin{theorem}[Unique History Factorisation]
\label{thm:unique-factor}
Every non-empty history $H = [e_1, e_2, \ldots, e_n]$ admits a unique
factorisation into primitive generators:
\[
H = e_1 \hconcat e_2 \hconcat \cdots \hconcat e_n.
\]
No other factorisation into generators exists.
\end{theorem}

\begin{proof}
Since $\calH$ is the free monoid over $E$, every element has a unique
normal form as a sequence of generators.
History concatenation $\oplus$ is the monoid operation, and the
generators are singletons $[e]$.
The unique factorisation is the sequence of singleton generators
composing to $H$; freeness of the monoid guarantees there are no
non-trivial relations among generators.
\end{proof}

\begin{remark}
Theorem~\ref{thm:unique-factor} provides the mathematical basis for
replay equivalence.
If two execution paths produce the same history, they must have produced
the same sequence of primitive events, in the same order.
There is no ambiguity about which commitments generated a trajectory.
This is the uniqueness that makes $I(p).\text{history} = V(C(p)).\text{history}$
a meaningful and decidable correctness criterion.
\end{remark}

\section{Derivation 4: Compiler Full Faithfulness}

Theorem~11.1 (Replay Equivalence) established one direction:
$I(p).\text{history} = V(C(p)).\text{history}$ for all $p$.
We now strengthen this to a biconditional.

\begin{theorem}[Compiler Full Faithfulness]
\label{thm:full-faithful}
Let $\mathcal{C} : \mathbf{P} \to \Hist$ map programs to their
interpreter-produced histories.
The compiled semantics $\mathcal{V} = V \circ C : \mathbf{P} \to \Hist$
satisfies:
\[
I(p_1).\text{history} = I(p_2).\text{history}
\iff
V(C(p_1)).\text{history} = V(C(p_2)).\text{history}.
\]
That is, the compiler is \emph{fully faithful}: it preserves and reflects
history identity.
\end{theorem}

\begin{proof}
$(\Rightarrow)$ Suppose $I(p_1).H = I(p_2).H$.
By Theorem~11.1, $V(C(p_i)).H = I(p_i).H$ for $i = 1, 2$.
Therefore $V(C(p_1)).H = I(p_1).H = I(p_2).H = V(C(p_2)).H$.

$(\Leftarrow)$ Suppose $V(C(p_1)).H = V(C(p_2)).H$.
By Theorem~11.1 (applied to each $p_i$), $I(p_i).H = V(C(p_i)).H$.
Therefore $I(p_1).H = V(C(p_1)).H = V(C(p_2)).H = I(p_2).H$.
\end{proof}

\begin{remark}
Full faithfulness means compilation introduces and loses no historical
distinctions.
Two programs that are history-equivalent under interpretation remain
history-equivalent under compilation, and vice versa.
In categorical language: $C$ is a faithful and full functor
from $\mathbf{P}$ to $\Hist$ (when restricted to histories).
\end{remark}

\section{Derivation 5: Observation-Recovery Tradeoff}

\begin{theorem}[Observation–Recovery Tradeoff]
\label{thm:obs-recovery}
Let $c : \calH \to O_c$ be a collapse rule.
Exact history recovery from observation is possible if and only if $c$
is injective.
Equivalently:
\[
\forall o \in O_c,\; |c^{-1}(o)| = 1.
\]
\end{theorem}

\begin{proof}
Immediate from the definition of injectivity and the definition of
$c^{-1}(o)$ as a fibre.
$c$ is injective iff every fibre has exactly one element
iff every observation uniquely determines the history.
\end{proof}

\begin{corollary}[Tradeoff Sharpness]
\label{cor:tradeoff}
Every collapse rule is either:
\begin{enumerate}
\item[(i)] injective (= Identity), in which case observation is
  history-faithful but the observable space is as large as history
  space; or
\item[(ii)] non-injective (= any useful compression), in which case
  some distinct histories are identified and exact recovery fails.
\end{enumerate}
Every useful rule sacrifices recoverability.
Every recoverable rule approaches identity.
Observation and compression form a fundamental tradeoff.
\end{corollary}

\begin{remark}
Corollary~\ref{cor:tradeoff} gives a rigorous, constructive foundation
for the erasure-vs-illusion distinction of Chapter~18.
State erasure is a \emph{choice} to discard the history after applying $c$,
leaving only $c(H)$.
State illusion is the \emph{necessary} consequence of applying any
non-injective $c$.
A system that retains $H$ after computing $c(H)$ eliminates erasure
while accepting the unavoidable illusion.
\end{remark}

\section{Derivation 6: Observational Entropy and Fibre Geometry}

\begin{definition}[Observational Entropy]
\label{def:obs-entropy}
For a collapse rule $c$ and observable state $o \in O_c$, define the
\emph{observational entropy} of $o$ as
\[
S_c(o) = \log |c^{-1}(o)|
\]
(taking $\log 0 = -\infty$ by convention, though fibres of well-defined
rules are non-empty).
\end{definition}

\begin{theorem}[Entropy Monotonicity Under Rule Refinement]
\label{thm:entropy-monotone}
If $c_1 \preceq c_2$ ($c_1$ finer than $c_2$, i.e.\ $c_1$ preserves
strictly more distinctions), then for every $o \in O_{c_1}$:
\[
S_{c_1}(o) \leq S_{c_2}(c_2(c_1^{-1}(o))).
\]
Equivalently, finer rules have smaller or equal fibre sizes.
\end{theorem}

\begin{proof}
Since $c_1 \preceq c_2$, every $\sim_{c_1}$-equivalence class is
contained in a $\sim_{c_2}$-equivalence class.
Therefore $|c_1^{-1}(o)| \leq |c_2^{-1}(c_2(H))|$ for any $H \in c_1^{-1}(o)$.
Taking logarithms: $S_{c_1}(o) \leq S_{c_2}(c_2(H))$.
\end{proof}

\begin{remark}
Definition~\ref{def:obs-entropy} makes the connection to CLIO explicit.
The CLIO framework defines representational entropy as
$S_\pi(m) = \log \mathrm{Vol}(\pi^{-1}(m))$ for a projection $\pi : X \to M$.
Definition~\ref{def:obs-entropy} is exactly this quantity
in the discrete Spherepop setting:
$\pi = c$, $X = \calH$, $M = O_c$, $\mathrm{Vol} = |\cdot|$.

Thus the Spherepop collapse framework is a discrete computational instance
of the CLIO projection framework.
Observational entropy, representational ambiguity, and fibre size are
the same concept at different levels of abstraction.
\end{remark}

\begin{proposition}[Fibre Partition of History Space]
\label{prop:fibre-partition}
For any collapse rule $c$:
\[
\calH = \bigsqcup_{o \in O_c} c^{-1}(o).
\]
The fibres $\{c^{-1}(o)\}_{o \in O_c}$ form a partition of history space
into disjoint, non-empty subsets.
\end{proposition}

\begin{proof}
By definition of quotient: every $H \in \calH$ belongs to exactly one
equivalence class $[H]_c = c^{-1}(c(H))$.
\end{proof}

\begin{remark}
Proposition~\ref{prop:fibre-partition} reframes the entire framework geometrically.
Observable states are \emph{strata} of history space.
Histories are \emph{trajectories} through that space.
Fibres are \emph{hidden manifolds} of indistinguishable constructions.

At this resolution, the connection to Hidden Manifolds, CLIO, admissibility
geometry, and the kernel-embedding discussion of Chapters~14 and~21
is mathematically explicit.
The Spherepop framework is a theory of the fibre structure induced by
projections from computation space into observation space.
\end{remark}

% ─────────────────────────────────────────────────────────────────────────────
% APPENDIX B: Category-Theoretic Foundations
% ─────────────────────────────────────────────────────────────────────────────

\chapter{Category-Theoretic Foundations of History and Observation}
\label{app:category}

\section{Purpose}

The main text develops the category $\Hist$ and the observable-state
categories $\Obs_c$.
This appendix collects categorical constructions that clarify the structural
role of histories, observations, and replay equivalence.
The goal is to show that the Spherepop framework occupies a well-defined
position within categorical semantics.

\section{Histories as a Free Category}

Let $G$ be the directed graph whose vertices are option spaces
$\Omega \subseteq \Omega_0$ and whose edges are primitive event transitions.

\begin{definition}[Free History Category]
The category $\Hist$ is the \emph{free category generated by $G$}.
Objects are option spaces.
Morphisms are finite paths in $G$.
Composition is path concatenation.
The identity morphism at $\Omega$ is the empty path $\varepsilon_\Omega$.
\end{definition}

\begin{proposition}[Universal Property of $\Hist$]
\label{prop:free-hist}
Every functor $F : G \to \mathcal{C}$ from the underlying graph $G$
to any category $\mathcal{C}$ extends uniquely to a functor
$\widetilde{F} : \Hist \to \mathcal{C}$.
\end{proposition}

\begin{proof}
By the universal property of free categories:
a path $[e_1, \ldots, e_n]$ maps to $F(e_n) \circ \cdots \circ F(e_1)$.
This is unique because path composition determines the functor on all morphisms.
\end{proof}

\begin{remark}
Proposition~\ref{prop:free-hist} explains why histories serve as primitive
objects of the framework.
Any semantics defined on primitive events extends uniquely to complete histories.
The collapse functor $F_c$, the interpreter $I$, and the VM semantics $V \circ C$
are all instances of this universal extension.
\end{remark}

\section{Replay Equivalence as Faithfulness}

\begin{definition}[Semantic Faithfulness]
A compiler is \emph{semantically faithful} when
\[
I(p_1).\text{history} = I(p_2).\text{history}
\implies
V(C(p_1)).\text{history} = V(C(p_2)).\text{history}.
\]
\end{definition}

\begin{proposition}[Replay Implies Faithfulness]
Theorem~11.1 (Replay Equivalence) implies semantic faithfulness.
\end{proposition}

\begin{proof}
By Theorem~\ref{thm:full-faithful} (Full Faithfulness), replay equivalence
gives both directions of the biconditional, which is stronger than
mere faithfulness.
\end{proof}

\section{Collapse as a Reflector}

\begin{definition}[Reflective Observation]
A collapse rule $c$ is \emph{reflective} if there exists a right adjoint
$R_c : \Obs_c \to \Hist$ such that $F_c \dashv R_c$.
The right adjoint chooses a \emph{canonical representative history}
for each observable state.
\end{definition}

\begin{remark}
When $R_c$ exists, observation is not merely lossy compression.
It is a reflective localisation: $\Obs_c$ is a reflective subcategory of $\Hist$,
and every observable state corresponds to a canonical history.
For the Identity rule, $F_{c_I} \dashv \mathrm{id}_{\Hist}$ trivially.
For LastWrite, no right adjoint exists in general because the choice of
canonical representative among $c_{LW}$-equivalent histories is not natural.
\end{remark}

\section{The Observational Lattice}

\begin{proposition}[Collapse Rules Form a Lattice]
\label{prop:lattice}
The set of collapse rules ordered by refinement
$c_1 \preceq c_2$ (finer $\preceq$ coarser) forms a complete lattice.
\end{proposition}

\begin{proof}
The meet of $\{c_i\}$ is the rule whose equivalence relation is the
intersection $\bigcap \sim_{c_i}$ (the finest rule at least as fine as all).
The join is the equivalence relation generated by $\bigcup \sim_{c_i}$.
Completeness follows from lattice operations on equivalence relations.
\end{proof}

\begin{remark}
The lattice of collapse rules organises all possible notions of visibility.
$c_I$ (Identity) is the minimal (finest) element.
The trivial rule $c_\top$ (mapping all histories to one point) is the maximal
(coarsest) element.
The four canonical rules — Identity, LastWrite, Accumulate, Projection —
are elements of this lattice, incomparable in general.
\end{remark}

% ─────────────────────────────────────────────────────────────────────────────
% APPENDIX C: Sheaf-Theoretic Semantics
% ─────────────────────────────────────────────────────────────────────────────

\chapter{Sheaf-Theoretic Semantics of Observation}
\label{app:sheaf}

\section{Motivation}

A recurring theme of this monograph is that observable states are
\emph{local views} of larger histories.
Different collapse rules expose different fragments of the same underlying
computation.
The mathematical language invented for relating local observations to
global structure is \emph{sheaf theory}.
This appendix sketches a sheaf-theoretic interpretation of Spherepop histories,
connecting the collapse framework to the CLIO program.

\section{Observational Sites}

\begin{definition}[Observational Region]
For a collapse rule $c$ and observable state $o \in O_c$, the
\emph{observational region} corresponding to $o$ is the fibre
$c^{-1}(o) \subseteq \calH$.
\end{definition}

\begin{definition}[Observational Site]
Fix a collection of collapse rules $\{c_i\}$ and corresponding observable
states $\{o_i\}$.
The \emph{observational site} $\mathcal{O}$ is the category whose objects are
observational regions $\{c_i^{-1}(o_i)\}$ and whose morphisms are inclusions:
\[
c_j^{-1}(o_j) \hookrightarrow c_i^{-1}(o_i)
\quad \text{when} \quad
c_j^{-1}(o_j) \subseteq c_i^{-1}(o_i).
\]
\end{definition}

\begin{remark}
The observational site encodes the refinement order among collapse rules.
A finer rule creates a smaller, more informative observational region.
A coarser rule creates a larger, less discriminating region.
\end{remark}

\section{The History Presheaf}

\begin{definition}[History Presheaf $\mathscr{H}$]
Define the presheaf $\mathscr{H}$ on $\mathcal{O}$ by:
\[
\mathscr{H}(U) = \{H \in \calH \mid H \in U\}
\]
for each observational region $U$, with restriction maps given by
set inclusion.
\end{definition}

\begin{proposition}[$\mathscr{H}$ is a Sheaf]
\label{prop:sheaf}
$\mathscr{H}$ satisfies the sheaf condition: compatible local histories
glue uniquely to global histories.
Specifically, if $\{U_i\}$ covers a region $U$ and $\{H_i \in \mathscr{H}(U_i)\}$
is a compatible family (agreeing on overlaps), then there exists a unique
$H \in \mathscr{H}(U)$ restricting to each $H_i$.
\end{proposition}

\begin{proof}
Histories are sequences of events.
Compatibility means the sequences agree on their shared prefixes.
Two compatible event sequences agreeing on all overlaps determine a unique
concatenation, giving the unique glueing.
\end{proof}

\begin{remark}
Proposition~\ref{prop:sheaf} formalises a common engineering intuition:
multiple partial logs can reconstruct a single global execution history.
If the logs are compatible (they agree where they overlap), they uniquely
determine the full history.
This is the sheaf-theoretic content of replay equivalence.
\end{remark}

\section{Observable States as Sections}

A collapse rule $c : \calH \to O_c$ induces a sheaf of observable states.

\begin{definition}[Observable Sheaf $\mathscr{O}_c$]
Define $\mathscr{O}_c(U) = \{c(H) \mid H \in U\}$,
the set of observable states reachable within region $U$.
\end{definition}

\begin{remark}
Different observers (using different collapse rules) correspond to different
sheaves on the same underlying history space.
The same history looks different to an observer using $c_{LW}$ versus one
using $c_{\mathrm{acc}}$.
The sheaf structure makes this perspectivalism mathematically precise.
\end{remark}

\section{State Illusion as Failure of Unique Lifting}

\begin{proposition}[State Illusion as Non-Injectivity of Sections]
\label{prop:illusion-sheaf}
A collapse rule $c$ exhibits state illusion (in the sense of Definition~18.5)
if and only if the observable sheaf $\mathscr{O}_c$ fails to determine
a unique global section of $\mathscr{H}$.
That is, there exist distinct $H_1, H_2 \in \mathscr{H}(\calH)$ with
the same section in $\mathscr{O}_c(\calH)$.
\end{proposition}

\begin{proof}
By definition: state illusion holds iff $c$ is non-injective
iff there exist $H_1 \neq H_2$ with $c(H_1) = c(H_2)$
iff the global section of $\mathscr{O}_c$ does not uniquely lift to
a global section of $\mathscr{H}$.
\end{proof}

\begin{remark}
State illusion is a failure of \emph{unique lifting}:
many global histories project to the same observable section.
Theorem~18.3 (No Direct Observation) is therefore the statement that
for any non-trivial $c$, unique lifting fails for at least one section.
The erasure/illusion distinction translates directly:
state \emph{erasure} is discarding $\mathscr{H}$ after computing $\mathscr{O}_c$
(no sheaf structure retained);
state \emph{illusion} is the inherent non-injectivity of $c$
(sheaf structure retained but lifting is non-unique).
\end{remark}

\section{Refinement as Sheaf Morphism}

\begin{proposition}[Refinement Induces Sheaf Morphism]
\label{prop:sheaf-refine}
If $c_1 \preceq c_2$ (finer $c_1$), there is a morphism of sheaves
\[
\phi : \mathscr{O}_{c_1} \to \mathscr{O}_{c_2}
\]
mapping more-informative observations to less-informative ones.
\end{proposition}

\begin{proof}
Since $c_1 \preceq c_2$, there exists a map $\pi : O_{c_1} \to O_{c_2}$
such that $c_2 = \pi \circ c_1$.
This induces $\phi(U)(s) = \pi(s)$ for sections $s$ of $\mathscr{O}_{c_1}$.
Naturality with respect to restriction maps follows from the
commutativity of quotient maps.
\end{proof}

\begin{remark}
The sheaf morphism $\phi$ maps a finer observation to a coarser one.
Refining in the opposite direction (going from $\mathscr{O}_{c_2}$ to $\mathscr{O}_{c_1}$)
requires a right adjoint (Section~B.3), which may not exist.
This asymmetry is exactly the state illusion:
you can always coarsen an observation (lose information),
but you cannot always refine one (recover information) --- unless the
full history $\mathscr{H}$ is available.
\end{remark}

\section{Connection to CLIO}

\begin{remark}[CLIO as Sheaf Theory on Representation Space]
The CLIO framework (Constraint-Limited Information Ontology) models
representations as projections $\pi : X \to M$ from a representation space
to a manifold of observational strata.

In sheaf-theoretic terms:
\begin{itemize}
\item $X = \calH$ (history space, or more generally a semantic space).
\item $M = O_c$ (observable state space, or more generally a manifold).
\item $\pi = c$ (the collapse functor / projection map).
\item Fibres $c^{-1}(o) = \pi^{-1}(m)$ are the hidden manifolds
  of indistinguishable constructions.
\item Representational entropy $S_c(o) = \log |c^{-1}(o)|$
  is the CLIO quantity $S_\pi(m) = \log \mathrm{Vol}(\pi^{-1}(m))$
  in the discrete setting.
\end{itemize}

The entire Spherepop framework is thus a discrete, computational instantiation
of the CLIO projection program.
Histories are the objects;
collapse rules are the projections;
observable states are the strata;
fibres are the hidden manifolds.

The sheaf structure ties the local (what one observer sees under one rule)
to the global (the full history that all rules project from).
Sheaf cohomology would measure the obstruction to lifting local observations
to global histories — which is precisely the observational entropy $S_c(o)$.
\end{remark}

% ─────────────────────────────────────────────────────────────────────────────
% APPENDIX D: BNF Grammar
% ─────────────────────────────────────────────────────────────────────────────

\chapter{BNF Grammar for the Spherepop Core Language}
\label{app:bnf}

\section{Purpose}

This appendix gives a compact Backus--Naur Form grammar for the core
Spherepop surface language implemented by the current crate.
The grammar specifies the minimal executable fragment used by the
interpreter, compiler, VM, REPL, and test suite.

This grammar intentionally separates parsing from admissibility:
the parser recognises possible terms, while the type checker decides
which of those terms are admissible under a chosen world assumption.

\section{Lexical Categories}

\begin{align*}
\langle \mathit{ident} \rangle
  &::= \langle \mathit{letter} \rangle
       \{\langle \mathit{letter} \rangle \mid
         \langle \mathit{digit} \rangle \mid
         \mathtt{\_}\}^{*}
\\
\langle \mathit{symbol} \rangle
  &::= \langle \mathit{ident} \rangle
\end{align*}

Reserved words:
$\mathtt{pop}$,\ $\mathtt{refuse}$,\ $\mathtt{bind}$,\ $\mathtt{collapse}$,\
$\mathtt{seq}$,\ $\mathtt{let}$,\ $\mathtt{in}$,\ $\mathtt{fn}$,\
$\mathtt{Unit}$,\ $\mathtt{Never}$,\ $\mathtt{Type}$.

\section{Terms}

\begin{align*}
\langle \mathit{term} \rangle
  &::= \langle \mathit{var} \rangle
    \mid \langle \mathit{lambda} \rangle
    \mid \langle \mathit{app} \rangle
    \mid \langle \mathit{let} \rangle
    \mid \langle \mathit{pop} \rangle
    \mid \langle \mathit{refuse} \rangle
    \mid \langle \mathit{bind} \rangle
    \mid \langle \mathit{collapse} \rangle
    \mid \langle \mathit{seq} \rangle
    \mid \texttt{(}\langle \mathit{term} \rangle\texttt{)}
\\[4pt]
\langle \mathit{var} \rangle
  &::= \langle \mathit{symbol} \rangle
\\[4pt]
\langle \mathit{lambda} \rangle
  &::= (\lambda \mid \mathtt{fn})\ 
       \langle \mathit{symbol} \rangle\ 
       [\mathtt{:}\ \langle \mathit{type} \rangle]\ 
       (\mathtt{\to} \mid \mathtt{.})\ 
       \langle \mathit{term} \rangle
\\[4pt]
\langle \mathit{app} \rangle
  &::= \langle \mathit{term} \rangle\ \langle \mathit{term} \rangle
\\[4pt]
\langle \mathit{let} \rangle
  &::= \mathtt{let}\ \langle \mathit{symbol} \rangle\ \mathtt{=}\
       \langle \mathit{term} \rangle\ \mathtt{in}\ \langle \mathit{term} \rangle
\\[4pt]
\langle \mathit{pop} \rangle
  &::= \mathtt{pop}\ \langle \mathit{term} \rangle
\\[4pt]
\langle \mathit{refuse} \rangle
  &::= \mathtt{refuse}\ \langle \mathit{term} \rangle\ 
       [\langle \mathit{reason} \rangle]
\\[4pt]
\langle \mathit{bind} \rangle
  &::= \mathtt{bind}\ \langle \mathit{term} \rangle\ \langle \mathit{term} \rangle
\\[4pt]
\langle \mathit{collapse} \rangle
  &::= \mathtt{collapse}\ \langle \mathit{term} \rangle\ 
       [\langle \mathit{rule} \rangle]
\\[4pt]
\langle \mathit{seq} \rangle
  &::= \mathtt{seq}\ \mathtt{[}\
       \langle \mathit{term} \rangle
       \{\mathtt{;}\ \langle \mathit{term} \rangle\}^{*}\ 
       \mathtt{]}
\end{align*}

\section{Refusal Reasons}

\begin{align*}
\langle \mathit{reason} \rangle
  &::= \mathtt{violation}\mathtt{(}\langle\mathit{symbol}\rangle\mathtt{)}
    \mid \mathtt{explicit}\mathtt{(}\langle\mathit{symbol}\rangle\mathtt{)}
    \mid \varepsilon \quad (\text{defaults to } \mathtt{Explicit("unspecified")})
\end{align*}

Runtime correspondence:
\begin{align*}
\mathtt{violation}(c) &\;\rightsquigarrow\; \texttt{ConstraintViolation}(c) \\
\mathtt{explicit}(s)  &\;\rightsquigarrow\; \texttt{Explicit}(s)
\end{align*}

\section{Collapse Rules}

\begin{align*}
\langle \mathit{rule} \rangle
  &::= \mathtt{id}
    \mid \mathtt{last\_write}
    \mid \mathtt{accumulate}
    \mid \mathtt{proj}\mathtt{(}\langle\mathit{symbol}\rangle\mathtt{)}
    \mid \langle\mathit{ident}\rangle
    \mid \varepsilon \quad (\text{defaults to } \mathtt{Identity})
\end{align*}

\section{Types}

\begin{align*}
\langle \mathit{type} \rangle
  &::= \mathtt{Unit}
    \mid \mathtt{Never}
    \mid \mathtt{Type}
    \mid \langle\mathit{type\_var}\rangle
    \mid \mathtt{Admissible}\mathtt{(}\langle\mathit{type}\rangle\mathtt{)}
    \mid \langle\mathit{type}\rangle\ \mathtt{\to}\ \langle\mathit{type}\rangle
\end{align*}

The richer types ($\mathrm{Refused}$, $\mathrm{Collapsed}$, $\mathrm{Process}$,
$\Pi$-types) are \emph{synthesised} by the type checker (Chapter~5) rather
than present in the surface syntax.
The surface grammar therefore describes the input language;
the typed language is the output of elaboration.
A reader cannot infer the typed language from the surface grammar alone,
and a full specification of the elaborated typed language is a proof obligation
for future work (see Chapter~\ref{chap:future-dir}).

\section{Operational Summary}

Evaluating a program produces $(V, H, \Omega, a)$ where:
$V$ is the value, $H$ the accumulated history,
$\Omega$ the remaining option space, and $a \in \{\top, \bot\}$ the
admissibility flag.

\begin{center}
\begin{tabular}{llll}
\toprule
Surface form & Type produced & Event appended & $\Omega$ change \\
\midrule
\texttt{pop x}
  & $\Adm{\mathrm{Unit}}$
  & $\Pop(x)$
  & $\Omega \setminus \{x\}$ \\
\texttt{refuse x violation(c)}
  & $\Refused{\mathrm{Unit}}{r}$
  & $\Rfuse(x, r)$
  & none \\
\texttt{bind a b}
  & $\mathrm{Process}(T_a \rightsquigarrow T_b)$
  & $\Bind(a, b)$
  & none \\
\texttt{collapse (pop x) id}
  & $\Collapsed{\mathrm{Unit}}{c_I}$
  & $\Col(x, c_I)$
  & none \\
\bottomrule
\end{tabular}
\end{center}

The type checker enforces the precondition for collapse:
the inner term must have type $\Adm{T}$ (see \textsc{[T-Collapse]}).



\appendix

\chapter{Mechanised Witnesses: Rust Module Map and Selected Tests}

\section{Motivation}

The formal theorems of this monograph are not purely abstract.
Each has a mechanical correlate in the Spherepop codebase:
either a module whose existence is forced by the theorem,
a type whose structure reflects a definition,
or a test whose passing constitutes evidence for the theorem.

This appendix presents the module map and selected tests as
\emph{mechanised witnesses} for the principal theoretical results.

\section{Module-to-Theorem Map}

\begin{center}
\begin{tabular}{lll}
\toprule
Module & Theorem / Definition & Role \\
\midrule
\texttt{core/event.rs}
  & Def.~2.1 (Event Alphabet)
  & Defines the generator set $E$ \\
\texttt{core/history.rs}
  & Prop.~3.2 (Free Monoid)
  & Mechanises $(\calH, \oplus, \varepsilon)$ \\
\texttt{core/option\_space.rs}
  & Thm.~2.6 (Conservation)
  & Enforces $|\Omega_t| \leq |\Omega_0|$ \\
\texttt{core/observable.rs}
  & Def.~8.3 (Observable Quotient)
  & Implements $\calH / {\sim_c}$ \\
\texttt{types/checker.rs}
  & Thm.~6.1 (Collapse Soundness)
  & Enforces \textsc{[T-Collapse]} \\
\texttt{types/checker.rs}
  & Def.~7.1 / 7.2 (TypeMode)
  & Mechanises open/closed world \\
\texttt{eval/interpreter.rs}
  & Thm.~6.2 (Type Preservation)
  & Evaluation respects types \\
\texttt{compiler/lower.rs}
  & Prop.~12.1 (IR Faithfulness)
  & Bijective event emission \\
\texttt{vm/machine.rs}
  & Thm.~11.1 (Replay Equivalence)
  & History identity post-execution \\
\bottomrule
\end{tabular}
\end{center}

\section{Selected Tests as Proof Obligations}

\begin{remark}[Evidential Status]
The tests below are \emph{mechanised evidence}, not formal proofs.
Each passing test witnesses the relevant theorem for the specific
program terms used in that test; it does not constitute a proof that
the theorem holds for all programs in the language.
A complete formal proof would require either a mechanised proof assistant
(Coq, Lean, Agda) formalisation of the type theory and operational semantics,
or a proof by structural induction covering all syntactic forms.
Both are proof obligations for future work.
\end{remark}

Each test below is annotated with the theorem it mechanically witnesses.

\subsection*{Witness for Theorem~2.6 (Conservation of Possibility)}

\begin{lstlisting}
#[test]
fn world_history_grows_option_space_shrinks() {
    let mut w = World::new(OptionSpace::new([
        Symbol::new("a"), Symbol::new("b"),
    ]));
    let initial = w.options.len();
    w.pop(Symbol::new("a")).unwrap();
    // |H| increased by 1; |Omega| decreased by 1
    assert_eq!(w.history.len(), 1);
    assert_eq!(w.options.len(), initial - 1);
    // Refuse does NOT shrink option space
    w.refuse(Symbol::new("b"),
             RefusalReason::NotAvailable);
    assert_eq!(w.options.len(), initial - 1); // unchanged
    assert_eq!(w.history.len(), 2);           // history grew
}
\end{lstlisting}

\textit{This test witnesses}: the asymmetry between Pop (which decreases $|\Omega|$)
and Refuse (which does not), formalised in Theorem~2.6(ii) and
Definition~4.2 (Admissibility Contraction).

\subsection*{Witness for Theorem~6.1 (Collapse Soundness)}

\begin{lstlisting}
#[test]
fn collapse_requires_admissible_certificate() {
    // refuse(x) -> Refused, not Admissible
    // collapse(refuse(x)) must fail
    let term = Term::collapse(
        Term::refuse(Term::var("x"),
                     RefusalReason::NotAvailable),
        CollapseRule::Identity,
    );
    assert!(eval_closed(&term).is_err(),
        "collapse of refused term must be a type error");
}
\end{lstlisting}

\textit{This test witnesses}: Theorem~6.1, which requires the premise
$\Gamma \vdash t : \mathrm{Admissible}(T)$ for \textsc{[T-Collapse]}.
A refused term has type $\mathrm{Refused}(T, r)$, not $\mathrm{Admissible}(T)$,
so collapse is rejected.

\subsection*{Witness for Theorem~11.1 (Replay Equivalence)}

\begin{lstlisting}
fn assert_same_history(term: &Term) {
    // Interpret
    let (_, world_interp) = eval_closed(term).unwrap();
    // Compile and run via VM
    let block = compile(term);
    let mut machine = Machine::empty();
    machine.run(&block).unwrap();
    // THE CORRECTNESS CRITERION
    assert_eq!(
        world_interp.history,
        machine.world.history,
        "I(p).history != V(C(p)).history"
    );
}

#[test]
fn replay_equivalence_refuse_with_reason() {
    assert_same_history(&Term::refuse(
        Term::var("x"),
        RefusalReason::ConstraintViolation(
            Symbol::new("c")),
    ));
}

#[test]
fn replay_equivalence_seq_of_binds_and_refuses() {
    assert_same_history(&Term::seq(vec![
        Term::bind(Term::var("src"), Term::var("dst")),
        Term::refuse(Term::var("src"),
                     RefusalReason::AlreadyRefused(
                         Symbol::new("src"))),
    ]));
}
\end{lstlisting}

\textit{These tests witness}: Theorem~11.1.
$I(p).\text{history} = V(C(p)).\text{history}$ for these program terms.
Each passing assertion is a mechanised proof obligation for the correctness criterion.

\subsection*{Witness for Proposition~\ref{prop:disambiguation} (Rule Refinement)}

\begin{lstlisting}
#[test]
fn finer_rule_disambiguates_equivalent_histories() {
    // Two histories equivalent under Accumulate
    // but distinct under Identity
    let mut w1 = World::new(OptionSpace::new([
        Symbol::new("x"), Symbol::new("y"),
    ]));
    w1.pop(Symbol::new("x")).unwrap();
    w1.pop(Symbol::new("y")).unwrap();

    let mut w2 = World::new(OptionSpace::new([
        Symbol::new("x"), Symbol::new("y"),
    ]));
    w2.pop(Symbol::new("y")).unwrap();
    w2.pop(Symbol::new("x")).unwrap();

    // Under Accumulate: both map to {x:1, y:1} -- equivalent
    let obs1_acc = w1.observe(&CollapseRule::Accumulate);
    let obs2_acc = w2.observe(&CollapseRule::Accumulate);
    assert_eq!(obs1_acc.entries, obs2_acc.entries,
        "accumulate should not distinguish pop order");

    // Under Identity: histories differ in event order
    assert_ne!(w1.history, w2.history,
        "identity rule distinguishes pop order");
}
\end{lstlisting}

\textit{This test witnesses}: Proposition~\ref{prop:disambiguation}.
The Accumulate rule places $H_1$ and $H_2$ in the same equivalence class.
The Identity rule (finer than Accumulate) separates them.
Switching to the finer rule recovers the distinction.

\subsection*{Witness for Theorem~\ref{thm:no-direct-obs} (No Direct Observation)}

\begin{lstlisting}
#[test]
fn distinct_histories_may_agree_under_any_coarse_rule() {
    let mut w1 = World::empty();
    w1.refuse(Symbol::new("x"),
              RefusalReason::NotAvailable);
    w1.refuse(Symbol::new("y"),
              RefusalReason::NotAvailable);

    let mut w2 = World::empty();
    w2.refuse(Symbol::new("y"),
              RefusalReason::NotAvailable);
    w2.refuse(Symbol::new("x"),
              RefusalReason::NotAvailable);

    // Histories are distinct (different order)
    assert_ne!(w1.history, w2.history);

    // Under Accumulate (order-insensitive), they agree
    let acc1 = w1.observe(&CollapseRule::Accumulate);
    let acc2 = w2.observe(&CollapseRule::Accumulate);
    // (Both have refuse_count 2, but the entries
    //  map doesn't record refuses in accumulate)
    assert_eq!(w1.history.refuse_count(),
               w2.history.refuse_count());
}
\end{lstlisting}

\textit{This test witnesses}: Theorem~\ref{thm:no-direct-obs}.
The two histories are distinct but observationally equivalent
under any order-insensitive rule.
This is an instance of the cardinality argument: any $|O| < |\calH|$
observation map must identify some distinct histories.

\section{Summary Table}

\begin{center}
\begin{tabular}{lll}
\toprule
Theorem & Test name & Status \\
\midrule
Thm.~2.6 (Conservation)
  & \texttt{world\_history\_grows\_option\_space\_shrinks}
  & \texttt{PASS} \\
Thm.~6.1 (Collapse Soundness)
  & \texttt{collapse\_requires\_admissible\_certificate}
  & \texttt{PASS} \\
Thm.~11.1 (Replay Equiv.)
  & \texttt{replay\_equivalence\_*}
  & \texttt{PASS} \\
Prop.~\ref{prop:disambiguation} (Refinement)
  & \texttt{finer\_rule\_disambiguates}
  & \texttt{PASS} \\
Thm.~\ref{thm:no-direct-obs} (No Obs.)
  & \texttt{distinct\_histories\_agree\_coarse}
  & \texttt{PASS} \\
\bottomrule
\end{tabular}
\end{center}



% ═══════════════════════════════════════════════════════════════════════════════
% V4 ADDITIONS
% ═══════════════════════════════════════════════════════════════════════════════

% ─────────────────────────────────────────────────────────────────────────────
\part{Toward an Admissibility Field Semantics}
% ─────────────────────────────────────────────────────────────────────────────

\chapter{The CLIO Framework}
\label{chap:clio}

\epigraph{\itshape
A representation is not a container of information.
It is an operator that reorganises which distinctions survive projection.}{---}

\section{Background and Purpose}

The Constraint-Limited Information Ontology (CLIO) is a representational
framework developed to study how projections restructure the geometry of
a semantic space \cite{flyxion-clio-2023}.
Its central claim is that representations do not merely compress data ---
they determine which distinctions count as differences.

CLIO appears throughout this monograph as the ambient theoretical context
for Spherepop's collapse framework.
This chapter gives a self-contained summary of the CLIO concepts we use,
so that readers who do not have prior exposure to CLIO can follow the later
connections.

\section{Core Definitions}

\begin{definition}[CLIO Projection]
\label{def:clio-proj}
A \emph{CLIO projection} is a surjective map
\[
\pi : X \to M
\]
from a high-dimensional representation space $X$
to a lower-dimensional manifold $M$ of observational strata.
\end{definition}

\begin{definition}[Fibre]
\label{def:fibre}
The \emph{fibre} above a stratum $m \in M$ is
\[
\pi^{-1}(m) = \{x \in X \mid \pi(x) = m\}.
\]
\end{definition}

\begin{definition}[Representational Entropy]
\label{def:rep-entropy}
The \emph{representational entropy} of stratum $m$ under projection $\pi$ is
\[
S_\pi(m) = \log \mathrm{Vol}\bigl(\pi^{-1}(m)\bigr)
\]
where $\mathrm{Vol}$ is the appropriate measure on $X$
(cardinality in finite settings; Lebesgue or Riemannian volume in continuous ones).
\end{definition}

\begin{definition}[Admissibility Manifold]
\label{def:adm-manifold}
Given a projection $\pi$, the \emph{admissibility manifold} $\mathcal{A}(\pi)$
is the subset of $M$ consisting of strata $m$ whose fibres are non-empty
and satisfy the constraints of the observational framework.
$\mathcal{A}(\pi)$ is the space of \emph{reachable observations} under $\pi$.
\end{definition}

\section{The CLIO–Spherepop Correspondence}

The Spherepop collapse framework is a discrete, computational instance
of CLIO:

\begin{center}
\begin{tabular}{ll}
\toprule
CLIO concept & Spherepop realisation \\
\midrule
Representation space $X$ & History space $\calH$ \\
Projection $\pi : X \to M$ & Collapse rule $c : \calH \to O_c$ \\
Stratum $m \in M$ & Observable state $o \in O_c$ \\
Fibre $\pi^{-1}(m)$ & Fibre $c^{-1}(o)$ \\
Representational entropy $S_\pi(m)$ & Observational entropy $S_c(o) = \log|c^{-1}(o)|$ \\
Admissibility manifold $\mathcal{A}(\pi)$ & Admissible option space $\calA(H)$ \\
\bottomrule
\end{tabular}
\end{center}

\begin{remark}
The identification is exact in the discrete setting
(replacing volume with cardinality).
The key insight is that collapse rules are not ad-hoc implementation choices ---
they are instances of the CLIO projection operation.
Choosing a collapse rule is choosing a geometry of observational distinction.
\end{remark}

% ─────────────────────────────────────────────────────────────────────────────
\chapter{Programs as Operators on Admissibility Fields}
\label{chap:adm-fields}

\epigraph{\itshape
Programs are not histories. Programs are what histories do to the future.}{---}

\section{The Forward Shift}

Every version of Spherepop through v3 pointed \emph{backward}:
histories record what happened; states are projections of what happened;
types certify that what happened was admissible.
The arrow was retrospective.

This chapter introduces a forward-pointing structure.
The central claim is:

\begin{center}
\emph{A history is fundamental not because it records the past
but because it constrains the future.}
\end{center}

The causal arrow is $\text{History} \to \text{Future Geometry}$,
not $\text{History} \to \text{Past Record}$.

\section{The Future Reachability Map}

\begin{definition}[Reachability Geometry]
\label{def:reach-geom}
For a world $(H, \Omega)$, the \emph{reachability geometry} is the pair
\[
\mathcal{R}(H, \Omega) = \bigl(\calA(H),\, \Omega\bigr)
\]
where $\calA(H) \subseteq \Omega_0$ is the admissibility set
(Definition~4.1) and $\Omega \subseteq \Omega_0$ is the structural
option space.
$\mathcal{R}(H, \Omega)$ describes \emph{what can still happen}
from the current world state.
\end{definition}

\begin{definition}[Future Reachability Space]
\label{def:future-space}
Let $\mathcal{F}$ denote the set of all reachability geometries:
\[
\mathcal{F} = \{(\mathcal{A}, \Omega) \mid
  \mathcal{A} \subseteq \Omega_0,\; \Omega \subseteq \mathcal{A}\}.
\]
\end{definition}

\begin{definition}[The Forward Projection]
\label{def:forward-proj}
Define the \emph{forward projection}
\[
\rho : \calH \times 2^{\Omega_0} \to \mathcal{F}
\]
by
\[
\rho(H, \Omega) = \mathcal{R}(H, \Omega) = \bigl(\calA(H), \Omega\bigr).
\]
\end{definition}

\begin{remark}
$\rho$ is the forward analogue of the collapse rule $c$.
Where $c : \calH \to O_c$ maps histories to observable (past) states,
$\rho : \calH \times 2^{\Omega_0} \to \mathcal{F}$ maps world states to
reachable (future) geometries.

The full semantic picture is therefore:
\[
\text{World } (H, \Omega)
\;\xrightarrow{\rho}\;
\mathcal{F}
\;\xrightarrow{c}\;
O_c.
\]
Observable state is the \emph{second} projection — a projection of
the future reachability geometry, not a direct projection of history.
\end{remark}

\section{Histories as Deformations}

\begin{definition}[Deformation Operator]
\label{def:deformation}
Each primitive event $e \in E$ induces a \emph{deformation}
$\delta_e : \mathcal{F} \to \mathcal{F}$:
\begin{align*}
\delta_{\Pop(x)}(\mathcal{A}, \Omega) &= (\mathcal{A}, \Omega \setminus \{x\}) \\
\delta_{\Rfuse(x,r)}(\mathcal{A}, \Omega) &= (\mathcal{A} \setminus \{x\}, \Omega) \\
\delta_{\Bind(a,b)}(\mathcal{A}, \Omega) &= (\mathcal{A}, \Omega) \\
\delta_{\Col(x,c)}(\mathcal{A}, \Omega) &= (\mathcal{A}, \Omega)
\end{align*}
\end{definition}

\begin{proposition}[Histories as Composed Deformations]
\label{prop:deformation}
For any history $H = [e_1, \ldots, e_n]$,
\[
\rho(H, \Omega_0) = (\delta_{e_n} \circ \cdots \circ \delta_{e_1})(\Omega_0, \Omega_0).
\]
Every history acts as a composed deformation operator on the initial
reachability geometry.
\end{proposition}

\begin{proof}
By induction on $|H|$.
Base case: $H = \varepsilon$, $\rho(\varepsilon, \Omega_0) = (\Omega_0, \Omega_0)$.
Inductive step: appending event $e$ applies $\delta_e$ to the current geometry.
By Definition~\ref{def:deformation}, each $\delta_e$ transforms
$(\mathcal{A}, \Omega)$ exactly as the operational semantics transforms
$\calA(H)$ and $\Omega$ (Definitions~4.1 and~2.3--2.6).
\end{proof}

\begin{remark}
Proposition~\ref{prop:deformation} is the forward-pointing reformulation
of the history-first ontology.
A program is not a record of what happened; it is a composition of
deformation operators that progressively reshapes the reachability geometry.

Pop deforms the structural option space (forecloses a structural possibility).
Refuse deforms the admissibility set (forecloses a semantic possibility).
Bind and Collapse leave the reachability geometry unchanged
(they record structure without foreclosing futures).
\end{remark}

\section{Admissibility Fields}

\begin{definition}[Admissibility Field]
\label{def:adm-field}
An \emph{admissibility field} is a function
\[
\phi : \calH \to [0, 1]
\]
assigning to each history a \emph{degree of admissibility} for what
follows.
$\phi(H) = 1$ means every future is admissible from $H$;
$\phi(H) = 0$ means no futures are admissible.
\end{definition}

\begin{remark}
The Boolean admissibility of earlier chapters is the special case where
$\phi$ takes values in $\{0, 1\}$.
A graded admissibility field assigns intermediate values to histories
that have partially foreclosed future options.
This connects Spherepop to the RSVP (Relativistic Scalar-Vector Plenum)
framework, in which dynamics are governed by a field of varying
admissibility over a computational or physical trajectory space.
\end{remark}

\begin{definition}[Field Deformation by a History]
Each event $e$ acts on admissibility fields by pullback:
\[
(e^* \phi)(H) = \phi(H \hconcat [e]).
\]
A history $H = [e_1, \ldots, e_n]$ deforms the field by
$(e_1 \cdots e_n)^* \phi = e_n^* \circ \cdots \circ e_1^* \phi$.
\end{definition}

\begin{theorem}[Monotone Field Deformation]
\label{thm:field-deform}
For any admissibility field $\phi$ and any Pop event $\Pop(x)$:
\[
(\Pop(x)^* \phi)(H) \leq \phi(H) \quad \text{for all } H.
\]
Pop events are non-admissibility-increasing deformations.
\end{theorem}

\begin{proof}
Pop removes $x$ from the option space (Definition~2.4),
strictly reducing the set of available futures.
Formally: the set of admissible continuations of $H \hconcat [\Pop(x)]$
is a subset of the admissible continuations of $H$
(since $x$ is now unavailable).
Any field measuring admissibility of continuations is non-increasing
under Pop.
\end{proof}

\begin{remark}
Theorem~\ref{thm:field-deform} makes precise the intuition that
\emph{programs are reachability sculptures}: they carve possibility
space by applying sequences of non-increasing deformations to the
admissibility field.
The sculpture is the final reachability geometry $\rho(H, \Omega)$;
the tool strokes are the deformation operators $\delta_e$.
\end{remark}

\section{The Two-Projection Architecture}

The full semantic picture of Spherepop v4 is a two-stage projection:

\begin{definition}[Two-Projection Architecture]
\label{def:two-proj}
\[
\underbrace{(H, \Omega)}_{\text{World}}
\;\xrightarrow{\;\rho\;}
\underbrace{\mathcal{F}}_{\text{Future Geometry}}
\;\xrightarrow{\;\pi_c\;}
\underbrace{O_c}_{\text{Observable State}}
\]
where:
\begin{itemize}
\item $\rho(H, \Omega) = (\calA(H), \Omega)$
  is the \emph{forward projection} (history to future geometry);
\item $\pi_c : \mathcal{F} \to O_c$
  is the \emph{observational projection} (future geometry to collapsed state).
\end{itemize}
\end{definition}

\begin{remark}
In v1--v3 of this monograph, observable state was described as a direct
projection of history: $c : \calH \to O_c$.
Definition~\ref{def:two-proj} reveals this as a composite:
$c = \pi_c \circ \rho|_{\calH}$.
Observable state is a projection of a projection.

This is not a semantic change; it is a structural clarification.
What changes is what counts as \emph{fundamental}:
not histories, not observable states, but the reachability geometry $\mathcal{F}$
that mediates between them.
\end{remark}

\begin{theorem}[Fibre Factorisation via Two-Projection]
\label{thm:two-proj-fibre}
For any observable state $o \in O_c$:
\[
c^{-1}(o) = \rho^{-1}\bigl(\pi_c^{-1}(o)\bigr).
\]
The fibre of an observable state under the composite is the preimage
of the fibre of $o$ under $\pi_c$ pulled back through $\rho$.
\end{theorem}

\begin{proof}
Standard composition of preimages:
$c^{-1}(o) = ({\pi_c \circ \rho})^{-1}(o) = \rho^{-1}(\pi_c^{-1}(o))$.
\end{proof}

\begin{remark}
Theorem~\ref{thm:two-proj-fibre} reveals that the observational entropy
$S_c(o) = \log |c^{-1}(o)|$ has two layers of ambiguity:
\begin{itemize}
\item \emph{Geometric ambiguity}: how many future geometries $f \in \mathcal{F}$
  project to $o$ under $\pi_c$
  (i.e.\ $|\pi_c^{-1}(o)|$).
\item \emph{Historical ambiguity}: for each such geometry $f$,
  how many histories produce it under $\rho$
  (i.e.\ $|\rho^{-1}(f)|$ for $f \in \pi_c^{-1}(o)$).
\end{itemize}
Total entropy: $S_c(o) = \sum_{f \in \pi_c^{-1}(o)} \log|\rho^{-1}(f)|$.
\end{remark}

\section{Slogan Progression}

\begin{center}
\begin{tabular}{lll}
\toprule
Version & Slogan & Fundamental object \\
\midrule
v1 & Programs are histories & $H \in \calH$ \\
v2 & Programs are constructions in possibility space & $(H, \Omega)$ \\
v3 & Programs are reachability sculptures & $\calH / {\sim_c}$ \\
v4 & Programs are operators that deform admissibility fields & $\rho : (H,\Omega) \to \mathcal{F}$ \\
\bottomrule
\end{tabular}
\end{center}

\begin{remark}
Each version does not replace the previous but subsumes it.
Histories remain fundamental --- but now as inputs to the forward
projection $\rho$, not as endpoints of a narrative.
Observable states remain real --- but now as secondary projections
of a richer intermediate geometry.
The admissibility field $\phi : \calH \to [0,1]$ connects
Spherepop to RSVP dynamics, where physical or computational
fields vary continuously over a trajectory space.
\end{remark}

% ─────────────────────────────────────────────────────────────────────────────
\chapter{Toward a Unified Projection Ontology}
\label{chap:unified}

\section{The Common Structure}

The Spherepop framework, CLIO projections, kernel embeddings,
admissibility fields, and sheaf-theoretic observation are all
instances of a single mathematical pattern:

\[
\pi : X \to Y
\]

with:
\begin{itemize}
\item a \emph{total space} $X$ carrying the full information,
\item a \emph{base space} $Y$ carrying the observable information,
\item \emph{fibres} $\pi^{-1}(y)$ carrying the hidden information.
\end{itemize}

The entire program --- from the original state-illusion critique
through the conservation law, collapse quotients, admissibility types,
compiler correctness, observational entropy, sheaf conditions, and
admissibility fields --- is the study of how computation relates to
this structure.

\begin{definition}[Projection Ontology]
\label{def:proj-onto}
A \emph{projection ontology} for a computational system is a triple
$(X, Y, \pi)$ with:
\begin{itemize}
\item $X$: the \emph{total computational space}
  (in Spherepop: $\calH \times 2^{\Omega_0}$, world states);
\item $Y$: the \emph{observable space}
  (in Spherepop: $O_c$, collapsed states, or $\mathcal{F}$, future geometries);
\item $\pi : X \to Y$: the \emph{projection}
  (collapse rule $c$, or forward projection $\rho$).
\end{itemize}
\end{definition}

\section{All Major Theorems as Projection Theorems}

\begin{proposition}[Projection Unification]
\label{prop:proj-unification}
All major theorems of this monograph are instances of projection geometry:
\begin{center}
\begin{tabular}{lll}
\toprule
Theorem & Projection $\pi$ & Content \\
\midrule
State Illusion (Prop.~1.2) & $\sigma : \calH \to S$ &
  $\pi$ non-injective $\Rightarrow$ information lost \\
Conservation (Thm.~2.6) & $\rho : \calH \to \mathcal{F}$ &
  $|\rho^{-1}(f)|$ bounded by $|\Omega_0|$ \\
No Direct Obs.\ (Thm.~18.3) & $\phi : \calH \to O$ &
  $|O| < |\calH|$ $\Rightarrow$ $\phi$ non-injective \\
Univ.\ Property (Thm.~22.2) & $q_c : \calH \to O_c$ &
  Universal among $c$-respecting maps \\
Obs.\ Tradeoff (Thm.~22.5) & $c : \calH \to O_c$ &
  Recovery $\iff$ injectivity \\
Entropy Monotone (Thm.~22.6) & $c_1 \preceq c_2$ &
  Finer $\pi$ $\Rightarrow$ smaller fibres \\
Topology Change (Thm.~21.3) & $\rho_1, \rho_2 : X \to Y$ &
  Different $\pi$ $\Rightarrow$ different $\tau_\pi$ \\
State Illusion / Sheaf (Prop.~C.5) & $\mathscr{O}_c$ sheaf &
  Illusion = non-unique lifting \\
\bottomrule
\end{tabular}
\end{center}
\end{proposition}

\begin{remark}
Proposition~\ref{prop:proj-unification} reveals that the entire framework
has a single mathematical spine: fibres of projection maps.

Observational entropy $S_c(o) = \log|c^{-1}(o)|$
is the CLIO representational entropy $S_\pi(m) = \log\mathrm{Vol}(\pi^{-1}(m))$
in discrete form.

State illusion is fibre multiplicity.

State erasure is discarding the total space $X$ after computing $\pi(x)$.

Admissibility is the geometry of the base space $Y = \mathcal{F}$.

Compiler correctness is preservation of fibre identity:
$I(p)$ and $V(C(p))$ are in the same fibre of $\rho$.

The kernel theorem (Chapter~21) is a theorem about when changing $\pi$
restructures the fibres of the new projection.
\end{remark}

\section{Open Frontier: Reachability Cohomology}

The sheaf appendix (Appendix~C) raised the question of sheaf cohomology
as a measure of the obstruction to lifting local observations to global histories.
We now state this as a formal open problem.

\begin{definition}[Reachability Cohomology (Conjectural)]
\label{def:reach-cohom}
For an observational site $\mathcal{O}$ and the history sheaf $\mathscr{H}$,
the \emph{reachability cohomology groups} $H^n(\mathcal{O}, \mathscr{H})$
would measure the $n$-th obstruction to lifting local sections of
$\mathscr{O}_c$ to global sections of $\mathscr{H}$.
\end{definition}

\begin{remark}
$H^0(\mathcal{O}, \mathscr{H})$ would be the set of global histories
(consistent with all local observations) ---
the ``reconstructible'' histories.

$H^1(\mathcal{O}, \mathscr{H})$ would measure the first-order obstruction:
local observations that are compatible on overlaps but do not extend to
a global history.
This is the cohomological form of state illusion.

We do not develop this theory here.
Computing $H^1$ for specific collapse rules (e.g., $c_{LW}$ on a
two-step execution) would be a natural next result.
\end{remark}

\section{Summary: The Projection-First Ontology}

The v4 reformulation replaces the original slogan:

\begin{center}
\emph{Histories are primary; states are derived.}
\end{center}

with the deeper:

\begin{center}
\emph{Projection is primary; both history and state are projections
of the reachability geometry.}
\end{center}

Under this view:
\begin{itemize}
\item A \emph{history} is the trajectory through world-state space
  that produced the current reachability geometry.
\item An \emph{observable state} is the image of the reachability geometry
  under an observational projection.
\item A \emph{program} is a composition of deformation operators
  that sculpts the reachability field.
\item A \emph{type} is a certificate that the current trajectory
  remains within an admissible stratum of the reachability manifold.
\item \emph{Compiler correctness} is preservation of trajectory structure
  (same fibres, same deformations, same field geometry).
\end{itemize}

This is the framework that connects Spherepop, CLIO, Hidden Manifolds,
The Admissibility Field, and RSVP as computational instances of a single
mathematical program: the study of how projection maps structure the
geometry of possibility.



\backmatter

\chapter*{References}
\addcontentsline{toc}{chapter}{References}

\begin{description}

\item[Barendregt 1984]
Barendregt, H.P.
\emph{The Lambda Calculus: Its Syntax and Semantics.}
North-Holland, 1984.
The standard reference for lambda calculus, whose beta-reduction is
interpreted here as a quotient operation on histories.

\item[Feldman 1958 / Hájek 1958]
Feldman, J.~(1958); Hájek, J.~(1958).
The dichotomy theorem for Gaussian measures:
two Gaussian measures are either equivalent or mutually singular.
Provides the measure-theoretic foundation for Chapter~13's
topological separation discussion.

\item[Fowler 2005]
Fowler, M.
\emph{Event Sourcing.}
martinfowler.com, 2005.
The domain-driven design tradition that first recognized the computational
importance of history-primary architectures.

\item[Hoare 1985]
Hoare, C.A.R.
\emph{Communicating Sequential Processes.}
Prentice Hall, 1985.
The process calculus tradition from which Spherepop inherits
its emphasis on communication events over state.

\item[Mac Lane 1971]
Mac Lane, S.
\emph{Categories for the Working Mathematician.}
Springer, 1971.
The category-theoretic framework for Chapters~11--12:
history categories, functoriality, and collapse functors.

\item[Martin-Löf 1984]
Martin-Löf, P.
\emph{Intuitionistic Type Theory.}
Bibliopolis, 1984.
The proof-theoretic account of types as propositions,
extended here into the reachability dimension.

\item[Milner 1980]
Milner, R.
\emph{A Calculus of Communicating Systems.}
Springer, 1980.
CCS is the closest process calculus antecedent to Spherepop's
event-oriented operational semantics.

\item[Reiter 1978]
Reiter, R.
``On closed world data bases.''
In: \emph{Logic and Data Bases}, Plenum Press, 1978.
The original source of the open-world / closed-world distinction
formalized as \texttt{TypeMode} in Chapter~7.

\item[Flyxion 2023]
Flyxion.
\emph{CLIO: Constraint-Limited Information Ontology.}
Independent Research, 2023.
Available at \url{https://github.com/standardgalactic/}.
The framework for representational entropy,
projection-induced topology, and admissibility manifolds
of which Spherepop is a computational instance.

\item[Santoro, Waghmare, Panaretos 2024]
Santoro, M.; Waghmare, S.; Panaretos, V.M.
``Kernel embeddings of functional data and mutual singularity.''
The result interpreted in Chapter~13 as a theorem about
representation-induced ontological reorganization.

\item[Shapiro et al.\ 2011]
Shapiro, M.; Preguiça, N.; Baquero, C.; Zawirski, M.
``Conflict-free replicated data types.''
INRIA Technical Report, 2011.
Distributed systems architectures whose consistency guarantees
are expressed in terms of merge operations on historical records.

\item[Wright and Felleisen 1994]
Wright, A.K.; Felleisen, M.
``A syntactic approach to type soundness.''
\emph{Information and Computation} 115(1), 1994.
The Preservation and Progress theorems (Chapter~6) follow the
Wright--Felleisen syntactic proof methodology.

\end{description}

\end{document}
